\documentclass[11pt,a4paper]{article}
\usepackage[T1]{fontenc}
\usepackage{lmodern,microtype}
\usepackage[margin=25mm]{geometry}
\usepackage{amsmath,amssymb,amsthm,mathtools,booktabs,enumitem}
\usepackage[hidelinks]{hyperref}
\hypersetup{pdftitle={Quantum-rational sums and holonomic boundary certificates}}

\newcommand{\Q}{\mathbb Q}
\newcommand{\Op}{\mathcal O}
\newcommand{\Bhat}{\widehat B}
\newcommand{\Ahat}{\widehat A}
\newcommand{\Bres}{B^{\mathrm{res}}}
\newcommand{\Bfor}{B^{\mathrm{for}}}
\newcommand{\Bhol}{B^{\mathrm{hol}}}
\newcommand{\Fhol}{F^{\mathrm{hol}}}
\newcommand{\Dop}{D_{\Op}}
\newcommand{\Khat}{\widehat K}
\newcommand{\val}{\nu}
\newcommand{\Mat}{\operatorname{Mat}}
\newcommand{\eps}{\varepsilon}
\newcommand{\Lt}{\widetilde L}
\newcommand{\dt}{\widetilde d}
\newcommand{\ft}{\widetilde f}

\theoremstyle{plain}
\newtheorem{theorem}{Theorem}[section]
\newtheorem{proposition}[theorem]{Proposition}
\newtheorem{lemma}[theorem]{Lemma}
\newtheorem{corollary}[theorem]{Corollary}
\theoremstyle{definition}
\newtheorem{definition}[theorem]{Definition}
\newtheorem{example}[theorem]{Example}
\theoremstyle{remark}
\newtheorem{remark}[theorem]{Remark}

\setlength{\parindent}{0pt}\setlength{\parskip}{4pt plus 1pt}
\setlist{nosep,leftmargin=*}
\allowdisplaybreaks[1]
\clubpenalty=10000 \widowpenalty=10000

\title{Quantum-rational sums and holonomic boundary certificates}
\author{}\date{}

\begin{document}
\maketitle\vspace{-3em}

\begin{abstract}
Let $\Delta$ be a rational element of the quantum plane $YX=tXY$, expanded
as a $t$-adic series, and let $S(\Delta)\in\Q[[t]]$ be the sum of its
coefficients at $X=Y=1$. We certify $S(\Delta)=0$ by a \emph{boundary
identity} $\Delta=(1-X)a+b(1-Y)$. The canonical boundary $b$ always exists
as a formal series and always satisfies a linear $t$-difference equation; the
whole difficulty is to prove that it \emph{converges}. We therefore split a
certificate into two independent parts: an algebraic identity that pins down
$b$, and a finite proof that its coefficients tend to $0$. For the first part
we give operator certificates rational in the shift~$T$, found by solving a
linear $q$-difference equation given by the least annihilator of the marginal;
those without a pole at $T=0$ prove convergence as well, and give $S(\Delta)$
exactly. For infinite product
identities a pole at $T=0$ is forced. For the second part we show that decay
along a single arithmetic
progression suffices, and give two finite tests: uniform contraction, and a
channel test with \emph{dilation witnesses}. Euler's odd--distinct partition
identity is proved this way. Since only the sum matters, each block of the
input may be replaced by other systems with the same sum: its second marginal,
obtained by exchanging $X$ and $Y$, and the second marginal of a canonical
realization. For partition generating functions these count by largest part
and by number of parts. Searching over these readings proves all $79$
differences within four families of realizations with a common sum, and all
$813$ differences within $497$ lists of automata with a common growth series:
most by a rational certificate without pole, and the infinite product
identities by the channel test. A sum side, whose partial sums are not
$q$-hypergeometric, is handled by channels up to a remainder that tends to $0$:
the channels are hypergeometric elements of the module of the partial sums,
obtained from the adjoint recurrence. With this, all $3237$ differences within
a second set of $49$ lists are proved except the $29$ that involve three
automata whose series equal the others by Gauss's theta-function identity. A short Julia
implementation built on the PBW algebras of \textsc{Oscar} accompanies the
paper.
\end{abstract}

\section{Overview}\label{sec:overview}

\paragraph{The problem.} Write $K=\Q(t)$ and let $A=K\langle X,Y\rangle/(YX-tXY)$
be the quantum plane. A rational element $\Delta$ of $A$ that is regular at
$t=0$ has a $t$-adic expansion $\Delta=\sum_k t^kH_k(X,Y)$ in normally ordered
polynomials $H_k$ (all $X$ before all $Y$). Its \emph{sum} is
\[
  S(\Delta)=\sum_k t^kH_k(1,1)\in\Q[[t]].
\]
Many $q$-series identities have the form $S(\Delta)=0$ for a difference
$\Delta=h-h'$ of two rational generating functions;
Section~\ref{sec:euler} treats Euler's identity in this form.

\paragraph{Boundary certificates.} If $\Delta=(1-X)a+b(1-Y)$ with $a,b$ in the
$t$-adic completion, then $S(\Delta)=0$ by telescoping. Conversely, whenever
$S(\Delta)=0$ there is a unique such pair with $b$ free of $X$
(Proposition~\ref{prop:canonical}). The function $b=b(Y)$ is the generating
function of the partial sums of the marginal $C=\rho(\Delta)$, so
\[
  \text{$b$ converges}\iff S(\Delta)=0.
\]
A \emph{complete} certificate consists of
\begin{enumerate}
\item an \textbf{identity} that determines $b$ by finite data, and
\item a finite \textbf{convergence proof} that the coefficients of $b$ tend to $0$.
\end{enumerate}
Part~1 is always possible: $b$ is holonomic (Theorem~\ref{thm:holonomic}). Part~2
is the real content; no complete method for it is known, and we do not claim one.

\paragraph{Results.}
\begin{itemize}
\item \emph{Rational operator certificates} (Section~\ref{sec:polynomial}),
  with coefficients rational in $T$, come from the rational solutions of a
  $q$-difference equation (Abramov's method), best formed from the least
  annihilator of the marginal (Proposition~\ref{prop:factor}). Without a pole
  at $T=0$ they deliver both parts at once and compute $S(\Delta)$; their
  existence is decidable.
\item \emph{Poles at $T=0$} (Section~\ref{sec:laurent}) are inverse shifts:
  such certificates deliver Part~1 only, and the order of the pole is the
  valuation that the convergence proof must make up. For identities between
  infinite products the pole is forced (Proposition~\ref{prop:forced-pole}).
\item \emph{One progression suffices} (Corollary~\ref{cor:progression}): the
  partial sums are $t$-adically Cauchy at a rate read off from the input, so
  decay along one arithmetic progression proves $S(\Delta)=0$, and then
  yields an explicit linear decay rate for all coefficients of $b$.
\item \emph{Convergence tests} (Section~\ref{sec:convergence}): uniform matrix
  contraction, and a channel test based on decimation and dilation witnesses,
  up to a finite boundary factor. The latter proves Euler's identity
  (Section~\ref{sec:euler}).
\item \emph{Remainders} (Section~\ref{sec:remainders}): partial sums that are not
  a sum of $q$-hypergeometric terms may be one up to a regular combination of
  themselves that tends to $0$. The terms are hypergeometric elements of their
  module, dual to the $q$-hypergeometric solutions of the adjoint recurrence,
  and the remainder is found by linear algebra on Laurent coefficients. This
  proves identities with a sum side.
\item \emph{Readings} (Section~\ref{sec:readings}): a block of the marginal
  system may be replaced by its second marginal (exchange of $X$ and $Y$) or by
  the second marginal of a canonical realization of its marginal, without
  changing $S$. Searching over readings proves all $79$ differences within
  four families of realizations with a common sum (Section~\ref{sec:families}),
  all $813$ differences within $497$ lists of automata with a common growth
  series, and, with remainders, all differences within a second set of $49$
  lists except those that need Gauss's identity (Section~\ref{sec:automata}).
\end{itemize}

\begin{center}
\begin{tabular}{@{}lll@{}}
\toprule
Route & Part 1 (identity) & Part 2 (convergence)\\\midrule
Rational certificate, no pole at $T=0$ & decidable & automatic\\
Rational certificate, pole of order $s$ & decidable & must gain $s$ per index\\
Vector telescoper & bounded ansatz & contraction or channel test\\
\bottomrule
\end{tabular}
\end{center}

\section{Algebras and series}\label{sec:spaces}

\subsection{Completions}
Let $\Khat=\Q((t))$ with valuation $\val=\val_t$, $\val(0)=+\infty$. Put
\[
\begin{array}{lll}
 A_0=\Q[t]\langle X,Y\rangle/(YX-tXY), & A=K\otimes A_0, & \Ahat=\varprojlim_N A_0/t^NA_0,\\[2pt]
 B_0=\Q[t][Y], & B=K[Y], & \Bhat=\varprojlim_N B_0/t^NB_0=\Q[Y][[t]].
\end{array}
\]
Elements of $\Ahat$ are series $\sum_{k\ge0}t^kH_k(X,Y)$ with $H_k\in\Q[X,Y]$
written in normal order; multiplication is the quantum product. The algebras
$B_0,B,\Bhat$ are the $X$-free parts. (Completion is taken before inverting
$t$: since $t$ is a unit in $A$, the $t$-adic completion of $A$ itself is $0$.)

Two larger series spaces carry convergence and its failure:
\[
  \Bres=\Bhat[t^{-1}]\qquad\text{and}\qquad \Bfor=\Khat[[Y]].
\]
For $v=\sum_n v_nY^n\in\Bfor$,
\begin{equation}\label{eq:membership}
  v\in\Bres\iff \val(v_n)\to+\infty,\qquad
  v\in\Bhat\iff \val(v_n)\to+\infty\ \text{and}\ \val(v_n)\ge0\ \text{for all }n.
\end{equation}
Thus $B\subset\Bhat\subset\Bres\subset\Bfor$ and $K[[Y]]\subset\Bfor$, but
$K[[Y]]\not\subset\Bres$. Evaluation at $Y=1$, $v\mapsto\sum_nv_n$, is defined
on $\Bres$ (a $t$-adically convergent sum) and maps $\Bhat$ to $\Q[[t]]$; it is
not defined on $\Bfor$. \emph{Formal existence in $\Bfor$ and convergence in
$\Bres$ are different assertions}, and the paper is organized around this distinction.

\subsection{The operator algebra and holonomic series}
Let $Y$ act on $\Bfor$ by multiplication and $T$ by $Tv(Y)=v(tY)$. Since
$TY=tYT$, this is an action of
\[
  \Op=K\langle Y,T\rangle/(TY-tYT),\qquad
  \Bigl(\sum_j h_j(Y)T^j\Bigr)\cdot v=\sum_j h_j(Y)\,v(t^jY).
\]
We write elements of $\Op$ in normal order $Y^bT^a$.

\begin{definition}
A series $v\in\Bfor$ is \emph{holonomic} if $H\cdot v=0$ for some
$0\ne H\in\Op$; equivalently, the $K(Y)$-span of $\{T^j v\}_{j\ge0}$ in
$\Khat((Y))$ is finite-dimensional. Write $\Bhol$ for the set of holonomic
series.
\end{definition}

The coefficient field is $K(Y)$, not $\Khat(Y)$; this is $q$-holonomicity in
$Y$ with $q=t$ (cf.~\cite{GL}), not differential finiteness in $t$.

\begin{proposition}\label{prop:closure}
$\Bhol$ is a $\Khat$-subalgebra of $\Bfor$ containing $\Khat[Y]$. It is stable
under $\Op$, under $T^{-1}v=v(Y/t)$, and under multiplication by elements of
$K(Y)$ whenever the product lies in $\Bfor$.
\end{proposition}
\begin{proof}
Let $V_v$ be the shift span of $v$. The shifts of $v+w$ lie in $V_v+V_w$, and
those of $vw$ in the span of products of bases of $V_v$ and $V_w$. On the
finite-dimensional space $V_v$ the map $T$ is injective and semilinear for the
automorphism $a(Y)\mapsto a(tY)$ of $K(Y)$, hence bijective; so $T^{-1}v\in V_v$
and $V_{T^{-1}v}=V_v$. The other statements are immediate.
\end{proof}

$\Bhol$ is \emph{not} a field: Appendix~\ref{app:not-field} gives
$g\in\Bhol\cap\Bhat$ with $g^{-1}\in\Bhat\setminus\Bhol$. Its fraction field
$\Fhol=\operatorname{Frac}(\Bhol)\subset\Khat((Y))$ consists of ordinary
quotients of commuting series; field membership carries no difference equation.
Note also that $\Bhol$ imposes no convergence: the relevant classes are
$\Bhol\cap\Bres$ and $\Bhol\cap\Bhat$.

\subsection{Operator fractions are not series}\label{sec:op-fractions}
The skew field $\Dop=\operatorname{Frac}(\Op)$ consists of noncommuting
operator fractions. It does not act on $\Bfor$: the nonzero operator $T-1$
kills the series $1$, whereas every nonzero element of $\Dop$ is invertible.
For instance $(1-tT)\cdot v=1$ has the unique solution
$v=(1-t)^{-1}$, while $(1-tT)^{-1}\in\Dop$ is not a scalar; its geometric
expansion acts on $\Bhat$ and gives $(1-tT)^{-1}\cdot1=(1-t)^{-1}$.
We use $\Dop$ only for symbolic elimination, and every inverse applied to a
series is given by an explicitly convergent expansion such as
\eqref{eq:geometric} below. We write $\eps(H)=H\cdot1\in K[Y]$; it is linear
but not multiplicative ($\eps(TY)=tY\ne Y=\eps(T)\eps(Y)$).

\section{From a realization to the marginal system}\label{sec:input}

\subsection{Realizations and the marginal}
A \emph{right realization} of $\Delta$ is
\begin{equation}\label{eq:input}
  wQ=r,\qquad \Delta=we,\qquad Q\in\Mat_n(A),\ r\in A^{1\times n},\ e\in A^{n\times1},
\end{equation}
with $Q$ invertible over the fraction skew field of $A$. We always assume the
\emph{normalization}
\begin{equation}\label{eq:normalization}
  Q\equiv I\pmod t,\qquad r,e\ \text{regular at } t=0
\end{equation}
coefficientwise in the normal basis, and select $w=rQ^{-1}\in\Ahat^{1\times n}$
by geometric inversion. (Convergence of $\Delta$ alone does not imply convergence
of all coordinates $w_i$ of an arbitrary realization.)

The \emph{marginal} and the \emph{sum} are the linear maps
\[
  \rho\Bigl(\sum_k t^kH_k(X,Y)\Bigr)=\sum_kt^kH_k(1,Y)\in\Bhat,\qquad S(v)=\rho(v)(1)\in\Q[[t]].
\]
$\rho$ is not multiplicative: $\rho(YX)=tY\ne\rho(Y)\rho(X)$.

\begin{lemma}[intertwining]\label{lem:intertwining}
Let $\Phi:A\to\Op$ be the anti-isomorphism $\Phi(X^aY^b)=Y^bT^a$. For
$w\in\Ahat[t^{-1}]$ and $q\in A$,
\[
  \rho(wq)=\Phi(q)\cdot\rho(w).
\]
Moreover $\ker\rho|_{\Ahat}=(1-X)\Ahat$.
\end{lemma}
\begin{proof}
For $w=X^iY^j$, $q=X^aY^b$: $wq=t^{ja}X^{i+a}Y^{j+b}$, so
$\rho(wq)=t^{ja}Y^{j+b}=Y^bT^a\cdot Y^j$. The kernel statement is polynomial
division by $1-X$ in each $t$-degree.
\end{proof}

In other words, $A/(1-X)A\cong B$ as a right $A$-module, and $\Phi$ turns the
right action into the left $\Op$-action. Applying the lemma to
\eqref{eq:input} gives the \emph{compiled triple}
\begin{equation}\label{eq:triple}
  L_{ji}=\Phi(Q_{ij}),\qquad f_j=\rho(r_j)\in B,\qquad d_i=\Phi(e_i),
\end{equation}
and the \emph{marginal system}
\begin{equation}\label{eq:compiled}
  L\cdot m=f,\qquad C:=\rho(\Delta)=d\cdot m,\qquad m_i=\rho(w_i)\in\Bhat.
\end{equation}
By \eqref{eq:normalization}, $L=I-N$ where every coefficient of $N$ has positive
valuation, and the selected solution is
\begin{equation}\label{eq:geometric}
  m=\sum_{j\ge0}N^j\cdot f,
\end{equation}
the unique solution in $\Bhat^n$. All later objects are computed from the
triple $(L,f,d)$ alone.

\subsection{The canonical boundary}

\begin{proposition}\label{prop:canonical}
For $\Delta\in\Ahat$ there exist $a,b\in\Ahat$ with
\begin{equation}\label{eq:boundary}
  \Delta=(1-X)a+b(1-Y)
\end{equation}
if and only if $S(\Delta)=0$. In that case there is exactly one such pair
with $b\in\Bhat$, namely
\[
  b=\frac{C}{1-Y},\qquad a=(1-X)^{-1}\bigl(\Delta-C(Y)\bigr).
\]
\end{proposition}
\begin{proof}
Telescoping gives necessity. Write $C=\sum_kt^kC_k(Y)$. The quotient
$C/(1-Y)$ has polynomial $t$-coefficients iff $C_k(1)=0$ for all $k$, i.e.\
iff $S(\Delta)=0$; then $\Delta-C$ has zero marginal and is divisible by $1-X$
(Lemma~\ref{lem:intertwining}). For an arbitrary pair, applying $\rho$ replaces
$b$ by $\rho(b)\in\Bhat$ and division reconstructs $a$; uniqueness follows from
injectivity of multiplication by $1-X$ and by $1-Y$.
\end{proof}

Write $C=\sum_nc_nY^n$ and $\beta_n=\sum_{j\le n}c_j$. Since $1-Y$ is a unit in
$\Bfor$, the \emph{formal boundary} always exists:
\begin{equation}\label{eq:partial-sums}
  b=\frac{C}{1-Y}=\sum_n\beta_nY^n\in\Bfor,\qquad
  b\in\Bhat\iff\beta_n\to0\iff S(\Delta)=0.
\end{equation}
Under \eqref{eq:normalization} all $c_n,\beta_n$ lie in $K$ and are integral.

\subsection{An unconditional rate}\label{sec:rate}
Choose integers
\[
  F_0\ge\deg_Yf,\qquad E\ge\deg_Yd,\qquad D\ge\max\{1,\deg_YL\},\qquad\alpha=F_0+E.
\]
Each factor $N$ in \eqref{eq:geometric} raises $t$-order by at least one and
$Y$-degree by at most $D$, so
\begin{equation}\label{eq:support}
  \deg_Y[t^k]\,m_i\le F_0+Dk,\qquad \deg_Y[t^k]\,C\le\alpha+Dk.
\end{equation}
Hence $\val(c_n)\ge\lceil (n-\alpha)/D\rceil$, and with
\begin{equation}\label{eq:rate}
  k_\Delta(n)=\max\Bigl\{0,\Bigl\lceil\frac{n+1-\alpha}{D}\Bigr\rceil\Bigr\}
\end{equation}
we obtain:

\begin{proposition}\label{prop:rate}
For every normalized input and all $n'\ge n\ge0$,
\[
  \val(S(\Delta)-\beta_n)\ge k_\Delta(n),\qquad \val(\beta_{n'}-\beta_n)\ge k_\Delta(n).
\]
Consequently:
\begin{enumerate}
\item if $\val(\beta_n)<k_\Delta(n)$ for some $n$, then $S(\Delta)\ne0$, and the
  leading term of $\beta_n$ is the leading term of $S(\Delta)$;
\item $S(\Delta)=0$ if and only if $\val(\beta_n)\ge k_\Delta(n)$ for all $n$.
\end{enumerate}
\end{proposition}

\begin{corollary}[one progression suffices]\label{cor:progression}
If $\val(\beta_{n_j})\to\infty$ along some subsequence $n_j\to\infty$, then
$S(\Delta)=0$, $b\in\Bhat$, and $\val(\beta_n)\ge k_\Delta(n)$ for every $n$.
\end{corollary}
\begin{proof}
$S(\Delta)=\lim_n\beta_n=\lim_j\beta_{n_j}=0$; then apply Proposition~\ref{prop:rate}(2).
\end{proof}

So a convergence proof need only show decay along one arithmetic progression,
at any rate; the uniform linear rate $k_\Delta$ then comes for free. Note that
$k_\Delta$ does not by itself decide anything: checking item~2 for all $n$
\emph{is} the zero problem.

\section{The boundary is always holonomic}\label{sec:holonomic}

\begin{lemma}[elimination]\label{lem:elimination}
Let $G\in\Mat_s(\Op)$ be nonsingular over $\Dop$ with $T$-degree at most $D$, and
fix a column index $k$. There are $U\in\Op^{1\times s}$ and $0\ne H\in\Op$ with
$UG=H\,e_k^{\mathsf T}$ and $\deg_TU\le(s-1)D$.
\end{lemma}
\begin{proof}
Seek $U_j=\sum_{i\le v}u_{ji}(Y)T^i$ over $K(Y)$ annihilating all columns
except the $k$-th. There are $s(v+1)$ unknowns and at most $(s-1)(v+D+1)$ linear
equations over $K(Y)$; for $v=(s-1)D$ a nonzero solution exists. The surviving
entry $H$ is nonzero because $G$ is nonsingular. Clear a common left denominator.
\end{proof}

\begin{theorem}\label{thm:holonomic}
For every normalized input, $m_i,C\in\Bhol\cap\Bhat$ and $b\in\Bhol$, with
effectively computable annihilators. Moreover
$S(\Delta)=0\iff b\in\Bhol\cap\Bhat$. Thus holonomicity imposes no further
existence condition; it provides no convergence proof.
\end{theorem}
\begin{proof}
Adjoin the constant $u=1$, selected by $(T-1)u=0$, $u(0)=1$. Then
\[
  G\cdot\begin{pmatrix}m\\u\\b\end{pmatrix}=0,\qquad
  G=\begin{pmatrix}L&-f&0\\0&T-1&0\\-d&0&1-Y\end{pmatrix}
\]
with $G$ nonsingular. Lemma~\ref{lem:elimination} gives $UG=(0,\dots,0,H)$,
a finite identity in $\Op$; applying it to the selected solution gives
$H\cdot b=0$. Other columns give annihilators of the $m_i$, and
Proposition~\ref{prop:closure} gives $C=d\cdot m\in\Bhol$. The last statement
is \eqref{eq:partial-sums}.
\end{proof}

No inverse from $\Dop$ is applied to a series here, so there is no conflict
with Section~\ref{sec:op-fractions}. An annihilator alone does not select $b$:
the selected $m$, $u=1$ and $(1-Y)b=dm$ do.

\begin{example}
$\Delta=1$ gives $b=(1-Y)^{-1}\in\Bhol\setminus\Bres$, described finitely by
$\bigl((1-tY)T-(1-Y)\bigr)b=0$, $b(0)=1$. Every coefficient has valuation $0$:
a finite description is not a convergence proof.
\end{example}

The component $a$ has a similar description by replacing $\Phi$ with
$\Psi(X^aY^b)=x^aY^bT^a$ over $\Op[x]$; no separate search for $a$ is needed.

\subsection{Identically vanishing marginals}\label{sec:zero}
An important special case is $C\equiv0$. It occurs whenever $\Delta=h-h'$ for
two realizations of the same series, or more generally of series with the same
marginal. Then $b=0$ and there is nothing to converge, yet a boundary
certificate may still need a shift (Example~\ref{ex:two-realizations}). We
certify $C=0$ directly by an annihilator. Write $H=\sum_bY^bH_b(T)$ for
$H\in\Op$.

\begin{proposition}[zero certificate]\label{prop:zero}
Let $0\ne H\in\Op$ and $R\in\Op^{1\times n}$ satisfy
\begin{equation}\label{eq:zerocert}
  H\,d=R\,L,\qquad \eps(Rf)=0.
\end{equation}
Let $b_0$ be least with $H_{b_0}\ne0$, and
$E=\{k\ge0:H_{b_0}(t^k)=0\}$. Then $E$ is finite and computable, and
$C=0$ if and only if $c_k=0$ for all $k\in E$.
\end{proposition}
\begin{proof}
$H\cdot C=(Hd)\cdot m=R\cdot(L\cdot m)=R\cdot f=\eps(Rf)=0$. Taking the
coefficient of $Y^{k+b_0}$,
\[
  H_{b_0}(t^k)\,c_k=-\sum_{b>b_0}H_b(t^{k+b_0-b})\,c_{k+b_0-b}.
\]
By the unique-minimal-term argument of Section~\ref{sec:convergence}, the
polynomial $H_{b_0}(z)$ is nonzero at $z=t^k$ for all $k$ beyond an explicit
bound, so $E$ is finite. If $c_k=0$ for $k\in E$, induction on $k$ gives $c=0$.
\end{proof}

\begin{proposition}\label{prop:zero-exists}
For every normalized input, a pair $(H,R)$ satisfying \eqref{eq:zerocert}
exists. Consequently it is decidable whether $C\equiv0$; in particular,
equality of the marginals of two realizations is decidable.
\end{proposition}
\begin{proof}
The matrix $G'=\left(\begin{smallmatrix}L&-f\\0&T-1\end{smallmatrix}\right)$ from
the proof of Theorem~\ref{thm:holonomic} is invertible over $\Dop$, so
$(d,0)=VG'$ for some row $V$ over $\Dop$. A common left denominator $H\ne0$ of
its entries (Ore condition) gives $U=HV=(R,R_u)$ over $\Op$ with $UG'=H(d,0)$,
that is, $RL=Hd$ and $Rf=R_u(T-1)$; applying the last identity to $1$ gives
$\eps(Rf)=0$. For the decision procedure, search $(H,R)$ by increasing degree,
which terminates because a solution exists. Then compute $E$, and compute
$c_k$ for $k\in E$.
\end{proof}

The equations \eqref{eq:zerocert} are linear and homogeneous in the
coefficients of $H$ and $R$. With a degree ansatz they are a kernel computation
over $K$.

\begin{example}[two realizations]\label{ex:two-realizations}
For the distinct-parts series $h'-1$ of Section~\ref{sec:euler}, the PBW
realization gives the marginal system $L=1-tY-t^2YT$, $f=1$, $d=tY$. A two-state
automaton gives $L'=\left(\begin{smallmatrix}1-tY&-tY\\-tT&1\end{smallmatrix}\right)$,
$f'=(1,0)$, $d'=(tY,0)$. Their difference is the block system
$\operatorname{diag}(L,L')$, $(1,1,0)$, $(tY,-tY,0)$ with $C\equiv0$. It has no
polynomial certificate. A certificate with a pole of order~$1$ at $T=0$ exists,
but its output is a cancelling combination of divergent series beyond the
channel test.
Proposition~\ref{prop:zero} applies at once:
\[
  H=1-tY-tYT,\qquad R=(tY,\,-tY,\,-t^2Y^2),\qquad E=\emptyset.
\]
For example $H\cdot tY=tY-t^2Y^2-t^3Y^2T=tY\cdot L$.
\end{example}

\section{Rational operator certificates}\label{sec:polynomial}

\paragraph{Rational coefficients act on series.} Since $T\cdot Y^n=t^nY^n$, a
rational function $r\in K(T)$ acts on $v=\sum v_nY^n$ coefficientwise,
\[
  r(T)\cdot v=\sum_n r(t^n)\,v_n\,Y^n,
\]
at every $n$ where $r$ has no pole at $t^n$. By the unique-minimum argument of
Section~\ref{sec:convergence}, a nonzero polynomial vanishes at $z=t^n$ for
finitely many computable $n$ only. So operators $\sum_bY^br_b(T)$ with
coefficients in $K(T)$ act on all coefficients of large index. This is the sense
in which certificates may be rational. (Fractions of operators in general do
not act, Section~\ref{sec:op-fractions}; fractions in $T$ alone do.)

\begin{definition}
A \emph{rational certificate} is a pair of rows $P,R$ with entries in $K(T)[Y]$ with
\begin{equation}\label{eq:polycert}
  d=(1-Y)P+RL.
\end{equation}
It is \emph{polynomial} if $R$ and $P$ have entries in $\Op$.
\end{definition}

\begin{proposition}\label{prop:polycert}
Let $(P,R)$ be a rational certificate and $q_n=[Y^n](P\cdot m)$. Then
$c_n=q_n-q_{n-1}+[Y^n](R\cdot f)$ for all large $n$, so $\beta_n-q_n=\kappa$ is
eventually constant. If no entry of $P$ has a pole at $T=0$, then $q_n\to0$ and
\[
  S(\Delta)=\kappa .
\]
If moreover no coefficient has a pole at any $t^n$, then $\kappa=(R\cdot f)(1)$
and $b=P\cdot m+(R\cdot f)/(1-Y)$.
\end{proposition}
\begin{proof}
Applying \eqref{eq:polycert} to $m$ gives $C=(1-Y)P\cdot m+R\cdot f$ at every
index where the coefficients involved are defined, and $R\cdot f$ is a polynomial.
If $p\in K(T)$ has no pole at $0$, then $\val p(t^j)$ is eventually constant, so
$\val(q_n)\ge\min_b\val(m_{n-b})-\mathrm{const}\to\infty$.
\end{proof}

So a rational certificate without pole at $T=0$ is a complete proof of
$S(\Delta)=0$, and otherwise computes $S(\Delta)\in K$ exactly. Denominators
that do not vanish at $T=0$ are harmless: $1/q(t^n)$ has eventually constant
valuation.

\paragraph{Reduction.} Let $\Lambda$ be evaluation at $Y=1$ in normal order,
$\Lambda(\sum Y^br_b(T))=\sum r_b(T)$, a map $K(T)[Y]\to K(T)$ with kernel
$(1-Y)K(T)[Y]$. Every row decomposes as $R=R_0+(1-Y)V$ with $R_0$ free of $Y$,
and replacing $(R,P)$ by $(R_0,P+VL)$ preserves \eqref{eq:polycert}. Hence a
rational certificate exists iff
\begin{equation}\label{eq:reduced}
  \Lambda(RL)=\Lambda(d),\qquad R\in K(T)^{1\times n},
\end{equation}
and $P$ is recovered by exact left division by $1-Y$: if $\Lambda(E)=0$ for
$E=\sum Y^be_b(T)$, then $E=(1-Y)P$ with $P=\sum Y^bp_b$,
$p_b=\sum_{b'\le b}e_{b'}$. In coordinates $z\leftrightarrow T$, with
$L_{ji}=\sum\ell_{ji,ba}Y^bT^a$ and $R_j=R_j(z)$, \eqref{eq:reduced} reads
\begin{equation}\label{eq:dual}
  \sum_{j,a,b}\ell_{ji,ba}\,z^aR_j(t^bz)=\Lambda(d_i)(z)\qquad(1\le i\le n),
\end{equation}
because $\Lambda(r(T)Y^bT^a)=r(t^bz)z^a$: a nonsingular linear $t$-difference
system for $R$ (the transpose of $L$ under $Y\mapsto S$, $T\mapsto z$, $Sz=tzS$).

\paragraph{Rational solutions.} Consider a scalar equation
$\sum_{b=0}^Ba_b(z)\,y(t^bz)=r(z)$ with polynomial coefficients, $a_0,a_B\ne0$
(Lemma~\ref{lem:elimination} applied to \eqref{eq:dual} gives one for each
$R_j$). Its rational solutions are found as follows~\cite{Abramov}.
\begin{enumerate}
\item \emph{Denominator.} Let $u$ be an irreducible factor of the denominator,
  $u\ne z$, and consider the chain of factors $u(t^iz)$ of the denominator. Its
  first element divides $a_0(z)$ and its last divides $A(z)=a_B(t^{-B}z)$, since
  otherwise that factor could not cancel in the equation. So the denominator divides
  $z^kU(z)$, where $U$ collects, for $s$ from the largest dispersion down to~$0$,
  $d(z)d(t^{-1}z)\cdots d(t^{-s}z)$ with $d=\gcd(A(z),a_0(t^sz))$, removing $d$
  from $A$ and $d(t^{-s}z)$ from $a_0$ at each step.
\item \emph{Pole at $0$.} If $y\sim cz^{-k}$, the lowest term of the left side
  is $cz^{v-k}\chi_0(t^{-k})$, where $v=\min_b\operatorname{ord}_za_b$ and
  $\chi_0(W)=\sum_b[z^v]a_b\,W^b$. So $k=v-\operatorname{ord}_zr$ or
  $\chi_0(t^{-k})=0$: finitely many $k$.
\item \emph{Numerator.} $N=z^kUy$ is a polynomial solution of an equation of
  the same kind, of degree bounded by Lemma~\ref{lem:degree}. Its coefficients
  solve a linear system over $K$, and the condition ``no pole at $0$'' is linear
  in them.
\end{enumerate}

\begin{lemma}[degree bound]\label{lem:degree}
Let $\sum_{j=0}^va_j(z)R(t^jz)=g(z)$ with $a_j,g\in K[z]$, not all $a_j=0$. Put
$D=\max_j\deg a_j$, $\chi(W)=\sum_j[z^D]a_j\,W^j$, and let $E$ bound the
$t$-degrees of the coefficients of $\chi$ after clearing denominators. Every
polynomial solution satisfies
\[
  \deg R\le\max\{0,E,\deg g-D\}\qquad(\text{omit }\deg g\text{ if }g=0).
\]
\end{lemma}
\begin{proof}
If $\deg R=N$, the coefficient of $z^{N+D}$ on the left is $\chi(t^N)$ times the
leading coefficient of $R$. For $N>E$ the terms of $\chi(t^N)$ occupy disjoint
$t$-degree ranges, so $\chi(t^N)\ne0$; this forces $N+D\le\deg g$.
\end{proof}

\begin{theorem}\label{thm:decidable}
Existence of a rational certificate, and of one without pole at $T=0$, is
decidable; such a certificate yields $S(\Delta)$ exactly. No bound in $Y$ is
needed.
\end{theorem}

Polynomial solution bounds for $q$-difference systems are also treated
in~\cite{Middeke}. The implementation uses the scalar form of
\eqref{eq:reduced} given by the annihilator: if $H\cdot C=g$, the $1\times1$
identity $1=(1-Y)P+R(T)H$ gives $C=(1-Y)P\cdot C+R\cdot g$, and the condition is
\begin{equation}\label{eq:scalar-cert}
  \Lambda(R\,H)=\sum_bh_b(t^bz)\,R(t^bz)=1,\qquad H=\sum_bh_b(T)Y^b
\end{equation}
(coefficients on the left), a scalar equation of order $\deg_YH$. Here
$b=P\cdot C$ is an image of $C$ itself; the system form can also succeed when
$b$ is an image of the auxiliary series but not of $C$, and its polynomial case
is searched as well.

\paragraph{The least annihilator.} The scalar form depends on $H$, and the
annihilator obtained from the system need not be the best one. The operators
$A\in K(T)[Y]$ for which the coefficients of $A\cdot C$ vanish at all large
indices form a left ideal of the Ore algebra $K(T)[Y]$, which is principal; let
$A_C$ be a generator. A certificate
$b=P\cdot C+(\text{polynomial})/(1-Y)$ exists iff $1-(1-Y)P$ lies in this ideal,
that is, by the reduction above, iff \eqref{eq:scalar-cert} has a rational
solution for $H=A_C$. A solution for a multiple $H=QA_C$ gives one for $A_C$
(reduce $R\,Q$ modulo $(1-Y)$), but not conversely. The annihilator from the
system is least for the class of $d$, not for $C$, and $A_C$ is often smaller:
when $f$ excites only part of the system, and typically for a difference
$C=C_1-C_2$ of two series that share a component. If $C$ is a polynomial,
$A_C=1$. A candidate $A$ is easy to guess from $c_0,\dots,c_N$ by linear
algebra; it is then proved as follows.

\begin{proposition}[a smaller annihilator]\label{prop:factor}
Let $H\cdot C=g$ with $g$ a polynomial, and let $H=QA$ with $A\in\Op$ and
$Q\in K(T)[Y]$. Choose $D\in K[T]$ with $\widetilde Q=DQ\in\Op$, let $g_A$ be
the part of $A\cdot C$ of degree less than $n_0$, and let $E$ be the exceptional
set of $\widetilde Q$ (Proposition~\ref{prop:zero}). If $\widetilde Q\cdot g_A=D\cdot g$ and the
coefficients of $A\cdot C$ vanish at the indices $n_0\le n\le\max E$, then
$A\cdot C=g_A$.
\end{proposition}
\begin{proof}
The series $u=A\cdot C-g_A$ satisfies
$\widetilde Q\cdot u=D\cdot(H\cdot C)-\widetilde Q\cdot g_A=D\cdot g-D\cdot g=0$.
Its coefficients vanish below $n_0$ by the choice of $g_A$, and up to $\max E$
by assumption. The induction in the proof of Proposition~\ref{prop:zero} gives
$u=0$.
\end{proof}

\begin{example}[a rational denominator]\label{ex:incomplete}
$\Delta=(1-t)(1-tX)^{-1}(X-1)$ compiles to $L=1-tT$, $f=1-t$, $d=T-1$. Here
\eqref{eq:reduced} reads $R(z)(1-tz)=z-1$, whose solution $R=(z-1)/(1-tz)$ is
not polynomial (nor Laurent) but has no pole at $z=0$ or at any $z=t^n$: a
complete certificate, with $S(\Delta)=R(1)(1-t)=0$. (In fact $C=0$ and
$\Delta=(1-X)a$; Proposition~\ref{prop:zero} also applies, with $H=1-tT$,
$R=T-1$ and $E=\emptyset$.) For a nonzero marginal, the system
$L=\left(\begin{smallmatrix}1-tY&0\\-t&1-tT^2\end{smallmatrix}\right)$, $f=(1,0)$,
$d=(0,1-Y)$ has $\beta_n=t^{n+1}/(1-t^{2n+1})$, and the scalar
certificate~\eqref{eq:scalar-cert} gives $b=P\cdot C$ with
$P=(t-T^2)/\bigl((t-1)(1+T^2)\bigr)$.
\end{example}

\section{Poles at \texorpdfstring{$T=0$}{T=0} and shifted unknowns}\label{sec:laurent}

A denominator prime to $T$ is harmless (Proposition~\ref{prop:polycert}). The
factor $T$ is not: $T^{-1}v=v(Y/t)$ preserves $\Bfor$ and $\Bhol$ but not $\Bres$,
\[
  (1-tY)^{-1}\in\Bhat,\qquad T^{-1}(1-tY)^{-1}=(1-Y)^{-1}\notin\Bres.
\]
Fix $s\ge0$ and change unknowns:
\begin{equation}\label{eq:shift}
  \mu=T^{-s}m,\quad \Lt=T^{-s}LT^s,\quad \ft=T^{-s}\cdot f,\quad \dt=dT^s,
  \qquad\text{so}\qquad \Lt\mu=\ft,\quad C=\dt\mu.
\end{equation}
Since $T^{-s}YT^s=t^{-s}Y$, conjugation only rescales coefficients:
$\Lt$ is again polynomial. The $\mu_i$ lie in $\Bhol$ but need not converge.

\begin{proposition}\label{prop:laurent}
Rational certificates whose entries have poles of order at most $s$ at $T=0$
correspond bijectively, via $P_s=PT^s$, $R_s=RT^s$, to rational certificates
$\dt=(1-Y)P_s+R_s\Lt$ of the shifted system without pole at $T=0$. Either form
proves
\[
  b=P_s\cdot\mu+(R_s\cdot\ft)/(1-Y)\in\Bhol,\qquad (1-Y)b=C,
\]
i.e.\ it identifies the canonical formal boundary. It does \emph{not} prove
convergence.
\end{proposition}
\begin{proof}
Multiply \eqref{eq:polycert} on the right by $T^s$; left division by $1-Y$ does
not change the order at $T=0$.
\end{proof}

\paragraph{What the shift costs.} If $g=T^sb=\sum\gamma_nY^n$ then
$\gamma_n=t^{sn}\beta_n$. A convergent expression for $g$ only gives
$\val(\gamma_n)\to\infty$; one needs $\val(\gamma_n)\ge sn+k(n)$ with
$k(n)\to\infty$. Inverse shifts move difficulty from the identity to the
estimate; they do not remove it. For identities between infinite products the
pole at $T=0$ cannot be avoided:

\begin{proposition}[forced pole]\label{prop:forced-pole}
Suppose that along some progression $n=pm+r$ the auxiliary sequences (or $C$,
in the scalar form) are sums of $q$-hypergeometric terms whose valuations tend
to infinity, while the decomposition of $\beta$ into terms with pairwise
non-similar ratios (Section~\ref{sec:intrinsic}) contains a term of bounded
valuation whose ratio is not similar to $1$. Then every rational certificate
has a pole at $T=0$.
\end{proposition}
\begin{proof}
Without a pole, $\beta_{pm+r}-\kappa$ is a combination of finitely many shifts of
the auxiliary sequences, with coefficients $r(t^{m})$ of bounded valuation for
large $m$. A $q$-hypergeometric term times such a coefficient is a term with a
similar ratio, of valuation still tending to infinity. Grouping by similarity
classes and using the uniqueness of the decomposition, every term of $\beta-\kappa$
would have valuation tending to infinity; the constant $\kappa$ only affects the
class of ratio~$1$.
\end{proof}

For Euler's identity, $\beta_{2m}=(t;t^2)_m^{-1}-(-t;t)_{2m}$ is a difference of
two terms of valuation $0$, while the auxiliary sequences have terms of
valuation $2m$ (Section~\ref{sec:euler}). So every certificate has a pole at
$T=0$; the limit $\beta_n\to0$ is the infinite product identity, and its
finite certificate is the dilation witness of Section~\ref{sec:convergence}, a
rational identity $\theta(\zeta)\chi(\zeta)=\prod_i\chi(t^{\mu+i}\zeta^\lambda)$
between the ratios of the two products. The same holds for the $41$ channel
pairs of Section~\ref{sec:families}: in the readings used, the least pole at
$T=0$ of a rational certificate~\eqref{eq:scalar-cert} has order~$2$.

\section{Convergence certificates}\label{sec:convergence}

\subsection{Vector form of the output}
For a normalized system with $L=\sum_{j=0}^qY^jL_j(T)$, coefficient
extraction gives
\[
  L_0(t^n)m_n=f_n-\sum_{j=1}^{q}L_j(t^{n-j})m_{n-j},
\]
with $L_0(t^n)$ invertible over $K$. (The same holds for $\Lt$ and $\mu$.)
Stacking $m_n,\dots,m_{n-q_i+1}$ (coordinate $i$ needs only lag
$q_i=\max_j\deg_YL_{ji}$) gives, beyond a known index $n_0$,
\begin{equation}\label{eq:state}
  v_{n+1}=A(t^n)v_n,\qquad \beta_n=\pi(t^n)\,v_n,\qquad A,\pi\ \text{over }K(z).
\end{equation}
An operator certificate yields $\pi$ by coefficient extraction:
$[Y^n]\,cY^jT^a\mu_i=c\,t^{a(n-j)}\mu_{i,n-j}$.

Alternatively, $\pi$ can be found directly as a \emph{vector telescoper}:
a rational row with
\begin{equation}\label{eq:telescoper}
  \pi(tz)A(z)-\pi(z)=\lambda(z),\qquad \pi(t^{N})v_{N}=\beta_{N},
\end{equation}
where $c_{n+1}=\lambda(t^n)v_n$; induction gives $\beta_n=\pi(t^n)v_n$ for
$n\ge N$. A bounded ansatz for $\pi$ is a linear system over $K$. Neither
discovery route is a prerequisite for the other, and both feed the same
convergence tests.

\paragraph{Valuations at $z=t^n$.} For $0\ne P(z)=\sum_ip_iz^i\in K[z]$ with
least exponent $i_0$, the term $i_0$ is the unique minimum of
$\val(p_i)+in$ as soon as $n>\max_{i>i_0}(\val p_{i_0}-\val p_i)/(i-i_0)$;
from then on $\val P(t^n)=\val(p_{i_0})+i_0n$. Applied to numerators and
denominators, every valuation inequality below is a finite check.

\subsection{Uniform contraction}
Choose integers $e_i$, $D(z)=\operatorname{diag}(z^{e_i})$,
$A_D(z)=D(tz)A(z)D(z)^{-1}$ and $\pi_D=\pi D^{-1}$. Suppose that for some
$\ell\ge1$ and $N$ every entry of $A_D(t^{\ell-1}z)\cdots A_D(z)$ at $z=t^n$
has valuation $\ge1$ for $n\ge N$, and every entry of $\pi_D(t^n)$ has valuation
$\ge h$. Then for $n\ge N$
\begin{equation}\label{eq:contraction}
  \val(\beta_n)\ge h+\min_{N\le j<N+\ell}\val\bigl(D(t^j)v_j\bigr)
  +\Bigl\lfloor\frac{n-N}{\ell}\Bigr\rfloor.
\end{equation}
This forces decay for \emph{every} initial vector, so it fails whenever decay
depends on cancellation specific to the selected initial values
(Remark~\ref{rem:obstruction}).

\subsection{Channels and dilation witnesses}
Fix a period $p\ge1$ and a residue $r$, and suppose the $p$-step product
\[
  A^{[p]}(x)=A(t^{p-1}x)\cdots A(tx)A(x)
\]
is \emph{diagonal}. Then along $n=pm+r$ each coordinate evolves by itself.
Put $\zeta=t^m$, choose $m_0$ with $pm_0+r\ge n_0$, and define the
\emph{channels}
\[
  G_k(m)=\pi_k(t^{pm+r})\,v_{pm+r,k},\qquad
  g_k(\zeta)=\frac{\pi_k(t^{r+p}\zeta^p)}{\pi_k(t^r\zeta^p)}\,A^{[p]}_{kk}(t^r\zeta^p).
\]
Then $G_k(m+1)=g_k(t^m)G_k(m)$ and $\beta_{pm+r}=\sum_kG_k(m)$ for $m\ge m_0$.

\begin{proposition}[channel test]\label{prop:channels}
Let $\mathcal K=\{k:G_k(m_0)\ne0\}$.
\begin{enumerate}[label=(\alph*)]
\item \emph{Decay.} If for each $k\in\mathcal K$ we have $\val g_k(t^j)\ge0$
  for $j\ge m_0$ and $\val g_k(t^j)\ge1$ for $j\ge M$, then
  $\val(\beta_{pm+r})\ge m+\min_k\val G_k(m_0)-M$.
\item \emph{Dilation.} Let $\mathcal K=\{k_1,k_2\}$ with $\val g_{k_1}(t^j)\ge0$
  for $j\ge m_0$, and put $\theta=g_{k_2}/g_{k_1}$. Suppose
  \begin{enumerate}[label=(\roman*)]
  \item the virtual initial ratio is $-1$:
    $G_{k_2}(m_0)=-\prod_{j<m_0}\theta(t^j)\,G_{k_1}(m_0)$;
  \item there are $\lambda\ge2$ and $\chi\in K[\zeta]$ with $\chi(0)=1$, integral
    coefficients, $\chi(1)\ne0$, and
    \begin{equation}\label{eq:dilation}
      \prod_{i=0}^{\lambda-1}\chi(t^i\zeta^\lambda)=\theta(\zeta)\,\chi(\zeta).
    \end{equation}
  \end{enumerate}
  Write $\chi-1=\sum_{i\ge i_1}\chi_i\zeta^i$ and $c=\min_i\val\chi_i$. Then
  \[
    \beta_{pm+r}=G_{k_1}(m)\Bigl(1-\prod_{j=m}^{\lambda m-1}\chi(t^j)\Bigr),
    \qquad
    \val(\beta_{pm+r})\ge\val G_{k_1}(m_0)+c+i_1m.
  \]
\end{enumerate}
In either case Corollary~\ref{cor:progression} gives $S(\Delta)=0$.
\end{proposition}
\begin{proof}
(a) $\val G_k(m)\ge\val G_k(m_0)+\max(0,m-M)$.
(b) First, $\prod_{j<m}\theta(t^j)=\prod_{j=m}^{\lambda m-1}\chi(t^j)$ for all
$m\ge0$, by induction: passing from $m$ to $m+1$ multiplies the right side by
$\prod_i\chi(t^{\lambda m+i})/\chi(t^m)=\theta(t^m)$ by \eqref{eq:dilation}
at $\zeta=t^m$; here $\chi(t^m)\ne0$ since $\chi(1)\ne0$ and $\chi(t^m)\equiv1$
for $m\ge1$. By (i) and the channel recursions,
$G_{k_2}(m)/G_{k_1}(m)=-\prod_{j<m}\theta(t^j)$ for $m\ge m_0$, which gives the
formula for $\beta$. Each $\chi(t^j)-1$ has valuation $\ge i_1j+c$, so
$\val\bigl(1-\prod_{j=m}^{\lambda m-1}\chi(t^j)\bigr)\ge i_1m+c$, and
$\val G_{k_1}(m)\ge\val G_{k_1}(m_0)$.
\end{proof}

\paragraph{Finding the witness.} Write $\theta=u/v$ in lowest terms. Comparing
degrees in $v(\zeta)\prod_i\chi(t^i\zeta^\lambda)=u(\zeta)\chi(\zeta)$ shows
\[
  (\lambda^2-1)\deg\chi=\deg u-\deg v,
\]
and comparing constant terms requires $\theta(0)=1$. The coefficient of
$\zeta^j$ on the left involves only coefficients of $\chi$ of degree
$\le j/\lambda$, so $\chi$ is determined degree by degree from $\chi(0)=1$.
One computes it up to the predicted degree and checks \eqref{eq:dilation}
exactly. The identity says $\prod_{j\ge0}\theta(t^j)=1$, i.e.\ it is a
finite certificate for an infinite product identity obtained by $\lambda$-section
of the products $E_m=\prod_{j<m}\chi(t^j)$, namely
$\prod_{j<m}\theta(t^j)=E_{\lambda m}/E_m$. The same recursion finds witnesses
with an \emph{offset} $\mu\in\mathbb Z$, i.e.\ solutions of
\begin{equation}\label{eq:dilation-mu}
  \prod_{i=0}^{\lambda-1}\chi(t^{\mu+i}\zeta^\lambda)=\theta(\zeta)\,\chi(\zeta);
\end{equation}
if $\chi$ is a witness for $\theta$, then $\chi(t^s\zeta)$ is a witness with
offset $\mu=s(\lambda-1)$ for $\theta(t^s\zeta)$. So offsets arise when the
channels are shifted in the index (Section~\ref{sec:six}).

The channel test is a restricted certificate class: it needs a diagonal
decimation, at most two live channels, and a dilation witness. It is not a
hypothesis on the input and carries no completeness claim.

\subsection{Channels of the output sequence}\label{sec:intrinsic}
The test above reads channels off a state space, so it depends on the
realization. An automaton realizing the same series typically has redundant
states: for the distinct-parts automaton of Example~\ref{ex:two-realizations},
the second coordinate is $t^{n+1}$ times the first. Then no product $A^{[p]}$ is
diagonal, and the test fails although the series is the same. We now describe a
version of the test that depends only on the sequence $(\beta_n)$.

\paragraph{A recurrence for $\beta$.} As in the proof of
Proposition~\ref{prop:zero-exists} (without the condition on $f$), there are
$H\ne0$ and $R$ with $Hd=RL$. Then $H\cdot C=\eps(Rf)=:g\in K[Y]$, and hence
$H(1-Y)\cdot b=g$. Write $H(1-Y)=\sum_bh_b(T)\,Y^b$ with the coefficients on
the left. Taking coefficients gives
\begin{equation}\label{eq:beta-rec}
  \sum_b h_b(t^n)\,\beta_{n-b}=g_n\qquad(n\ge0),
\end{equation}
a linear recurrence for $\beta$ that is homogeneous for $n>\deg g$. In this
normal order
\[
  (h(T)Y^b)(\kappa(T)Y^c)=h(T)\,\kappa(t^{-b}T)\,Y^{b+c},
\]
so the unknown coefficients of $H$ and $R$ enter linearly. Thus an $H$ of least
$Y$-degree, which is the order of \eqref{eq:beta-rec}, is found by linear
algebra over $K(T)$; one exists of degree $\le nD$ by the count of
Lemma~\ref{lem:elimination}. This is the same computation as in
Section~\ref{sec:zero}: if $g=0$, Proposition~\ref{prop:zero} decides whether
$C\equiv0$.

\paragraph{Sections.} Fix a period $p$ and a residue $r$, and put
$x_m=\beta_{pm+r}$ and $\zeta=t^m$. Write \eqref{eq:beta-rec} in companion form
$v_{k+1}=A(t^k)v_k$, and put $B(\zeta)=A^{[p]}(t^r\zeta^p)$. The rows
$w_0=e_1$, $w_{j+1}(\zeta)=w_j(t\zeta)B(\zeta)$ satisfy
$x_{m+j}=w_j(t^m)v_{pm+r}$. The first linear dependence among them over
$K(\zeta)$ gives a proven recurrence
\begin{equation}\label{eq:section-rec}
  \sum_{i=0}^{k}q_i(t^m)\,x_{m+i}=0\qquad(m\ge m_0),\qquad q_i\in K[\zeta].
\end{equation}

\paragraph{Channels.} A \emph{$q$-hypergeometric term} is a sequence with
$G(m+1)=g(t^m)\,G(m)$ for a rational $g\in K(\zeta)$, its \emph{ratio}. The
$K$-rational ratios of all $q$-hypergeometric solutions of
\eqref{eq:section-rec} are computed by the $q$-analogue of Petkov\v{s}ek's
algorithm~\cite{APP}. Write $g=Z\cdot A(\zeta)/B(\zeta)\cdot C(t\zeta)/C(\zeta)$
with $A\mid q_0$ and $B\mid q_k(t^{1-k}\zeta)$. The lowest and highest
coefficients of the resulting equation determine $Z$ and $\deg C$, and $C$ is
then found by linear algebra. Each ratio returned is verified by substitution.

\begin{proposition}[decomposition]\label{prop:decomp}
Let $q_k(t^m)\ne0$ and $g_i(t^m)\notin\{0,\infty\}$ for $m\ge m_1\ge m_0$, let
$g_i$ be ratios of solutions of \eqref{eq:section-rec}, and put
$G_i(m)=\prod_{j=m_1}^{m-1}g_i(t^j)$. If $x_m=\sum_ia_iG_i(m)$ for the $k$
values $m_1\le m<m_1+k$, then this holds for all $m\ge m_1$.
\end{proposition}
\begin{proof}
Both sides solve \eqref{eq:section-rec} for $m\ge m_1$. Since
$q_k(t^m)\ne0$, each value is determined by the $k$ preceding ones.
\end{proof}

\begin{proposition}[convergence of channels]\label{prop:channels2}
Let $x_m=\sum_ia_iG_i(m)$ for $m\ge m_1$, with $\val g_i(t^j)\ge0$ for $j\ge m_1$
whenever $a_i\ne0$. Call channel $i$ \emph{decaying} if $\val g_i(t^j)\ge1$ for
all large $j$, and \emph{stationary} otherwise. Suppose the stationary channels
can be grouped into pairs $(k_1,k_2)$ such that
\[
  \theta=\frac{g_{k_2}}{g_{k_1}}=\theta_w(\zeta)\,\frac{\Gamma(t\zeta)}{\Gamma(\zeta)},
  \qquad \Gamma\in K(\zeta),\quad \Gamma(0)=1,\quad
  \Gamma(t^j)\notin\{0,\infty\}\ (j\ge m_1),
\]
where $\theta_w$ satisfies \eqref{eq:dilation-mu} for some $\lambda\ge2$, offset
$\mu$ with $\lambda m_1+\mu\ge m_1$, and $\chi$ as in
Proposition~\ref{prop:channels}(ii), and
\[
  a_{k_2}/a_{k_1}=-F(m_1)\,\Gamma(t^{m_1}),\qquad F(m)=\prod_{l=m}^{\lambda m+\mu-1}\chi(t^l).
\]
Then $\val(x_m)\ge sm+\kappa$ for explicit $s\ge1$ and $\kappa$, and
$S(\Delta)=0$.
\end{proposition}
\begin{proof}
A decaying channel has $\val(a_iG_i(m))\ge\val a_i+m-M_i$, as in
Proposition~\ref{prop:channels}(a). For a pair, \eqref{eq:dilation-mu} at
$\zeta=t^m$ gives $F(m+1)/F(m)=\prod_i\chi(t^{\lambda m+\mu+i})/\chi(t^m)=\theta_w(t^m)$,
so $G_{k_2}(m)/G_{k_1}(m)=\prod_{j=m_1}^{m-1}\theta(t^j)
=F(m)\Gamma(t^m)/\bigl(F(m_1)\Gamma(t^{m_1})\bigr)$ and
\[
  a_{k_1}G_{k_1}(m)+a_{k_2}G_{k_2}(m)=a_{k_1}G_{k_1}(m)\bigl(1-F(m)\,\Gamma(t^m)\bigr).
\]
Every factor of $F(m)$ is some $\chi(t^l)$ with $l\ge m$, so
$\val(1-F(m))\ge i_1m+c$ as in Proposition~\ref{prop:channels}(b); and since
$\Gamma(0)=1$, $\val(1-\Gamma(t^m))\ge\sigma m+\tau$ for large $m$, with
$\sigma\ge1$. This proves decay along one progression, and
Corollary~\ref{cor:progression} concludes.
\end{proof}

\paragraph{Shifted factors.} With $\Gamma=1$ the pairing is that of
Proposition~\ref{prop:channels}(b). The factor $\Gamma$ is needed when the two
products differ by a finite boundary factor. For instance,
$(-t;t)_{2m}$ and $1/(t^3;t^2)_m$ have $\theta=(1+t\zeta^2)(1+t^2\zeta^2)(1-t^3\zeta^2)$,
which has no witness, but
\[
  (1-t)\,(-t;t)_{2m}\,(t^3;t^2)_m=(1-t^{2m+1})\prod_{l=m+1}^{2m}(1-t^{2l}),
\]
so $\Gamma(\zeta)=1-t\zeta^2$, $\theta_w=(1-t\zeta^2)(1+t\zeta^2)(1+t^2\zeta^2)$
and $\chi=1-t^2\zeta^2$ with $\lambda=2$, $\mu=0$. In general, $\theta_w$ is
obtained by replacing each irreducible factor $f$ of $\theta$, normalized to
$f(0)=1$, by another member $f(t^s\zeta)$ of its class
$\{f(t^s\zeta)\}_{s\in\mathbb Z}$; the implementation tries the members whose lowest coefficient
has valuation in a window $[w,w+e)$, $e$ the exponent of that coefficient, for
$w=0,\dots,3$, and then $\Gamma$ is a product of shifts of the $f$.

\paragraph{Independence of the realization.} Call two ratios $g,g'$
\emph{similar} if $g'/g=h(t\zeta)/h(\zeta)$ for some $h\in K(\zeta)^\times$.
Terms with pairwise non-similar ratios are linearly independent over
$K(\zeta)$; this is the $q$-analogue of the classical statement for
hypergeometric terms~\cite{PWZ}. Hence a sequence has at most one decomposition
into $q$-hypergeometric terms with pairwise non-similar ratios. Moreover, if
$x=\sum_iG_i$ satisfies a recurrence $Q$, then applying $Q$ gives
$\sum_ir_iG_i=0$ with $r_i\in K(\zeta)$, so every $r_i=0$: each constituent
solves every recurrence of $x$. So when a section of $\beta$ is a sum of
$q$-hypergeometric terms, its constituents are found from the recurrence
\eqref{eq:section-rec} of \emph{any} realization. The outcome of
Propositions~\ref{prop:decomp} and~\ref{prop:channels2} depends only on the
sequence $\beta$, within the search bounds (the periods $p$ tried, and the cap
on $\deg C$). This route also subsumes the shifted certificates of
Section~\ref{sec:laurent} for the purpose of proving convergence.

\paragraph{Merging similar channels.} The ratios returned by the
$q$-Petkov\v{s}ek step may be similar. Then the corresponding terms are
proportional with rational factors, the coefficients $a_i$ of
Proposition~\ref{prop:decomp} are not unique, and the individual terms can be
stationary although their sum decays. So each similarity class is first
replaced by one term: if $g_j=g_i\,h_j(t\zeta)/h_j(\zeta)$, then
$G_j(m)=G_i(m)\,h_j(t^m)/h_j(t^{m_1})$ and
\[
  \sum_ja_jG_j(m)=P(t^m)\,G_i(m),\qquad
  P(\zeta)=\sum_ja_j\,\frac{h_j(\zeta)}{h_j(t^{m_1})},
\]
a single term with ratio $g_i\,P(t\zeta)/P(\zeta)$; classes with $P=0$ disappear.
A solution $h$ of $h(t\zeta)/h(\zeta)=\rho$ is found from the factorization of
$\rho$. Its irreducible factors other than $\zeta$ fall into classes
$\{f(t^s\zeta)\}_{s\in\mathbb Z}$; $h$ exists iff $\zeta$ does not occur, the
exponents in each class sum to zero and the remaining constant is a power
$t^k$, and then $h=\zeta^k\prod f(t^s\zeta)^{-w_s}$, where $w_s$ is the sum of the
exponents of the class at shifts $\le s$.

\subsection{Channels up to a convergent remainder}\label{sec:remainders}
Proposition~\ref{prop:channels2} needs a section of $\beta$ that is a sum of
$q$-hypergeometric terms. This fails for a \emph{sum side}: a series such as
$\sum_nt^{\binom{n+1}2}/(t;t)_n$, whose partial sums are not
$q$-hypergeometric, while its product form $(-t;t)_\infty$ appears only in
another realization (Section~\ref{sec:automata} has examples). Such a side is
still a sum of $q$-hypergeometric terms up to a remainder that tends to $0$, and
the remainder is a combination of the sequence itself with coefficients
regular at $\zeta=0$.

\begin{proposition}[remainders]\label{prop:remainder}
Let $x_m\to S$ $t$-adically, and let $\pi_0,\dots,\pi_K\in K(\zeta)$ be regular
at $\zeta=0$ with $\sum_j\pi_j(0)=0$. Then $\sum_j\pi_j(t^m)\,x_{m+j}\to0$.
\end{proposition}
\begin{proof}
The sum is $\sum_j\pi_j(t^m)(x_{m+j}-S)+S\sum_j\pi_j(t^m)$. The valuations of
the $\pi_j(t^m)$ are bounded below, and $\sum_j\pi_j(t^m)\to\sum_j\pi_j(0)=0$.
\end{proof}
This applies to $x_m=\beta^{(b)}_{pm+r}$, the partial sums of one diagonal block
$b$ of the system: $x_m\to S^{(b)}$ by Proposition~\ref{prop:rate}, whatever the
value $S^{(b)}$ is.

\paragraph{Hypergeometric elements.} Let $x$ solve the section recurrence
\eqref{eq:section-rec},
\[
  (\mathcal Qx)_m=\sum_{i=0}^kq_i(t^m)\,x_{m+i}=0\qquad(m\ge m_0).
\]
A \emph{hypergeometric element} is a combination
$w_m=\sum_{j=1}^k\ell_j(t^m)x_{m+j}$ with $\ell_j\in K(\zeta)$ that is a
$q$-hypergeometric term. These are dual to the $q$-hypergeometric solutions of
the adjoint recurrence
\[
  (\mathcal Q^*y)_m=\sum_{i=0}^kq_i(t^{m-i})\,y_{m-i}=0.
\]
Indeed, the concomitant
\[
  W_m=\sum_{i=1}^k\sum_{l=m-i+1}^{m}q_i(t^l)\,y_l\,x_{l+i}
\]
satisfies $W_m-W_{m-1}=y_m(\mathcal Qx)_m-x_m(\mathcal Q^*y)_m$, so it is
constant once both recurrences hold. If $y_{m+1}=\gamma(t^m)y_m$, then
$w_m=W_m/y_m$ has ratio $1/\gamma$, and
\[
  \ell_j(\zeta)=\sum_{i=j}^kq_i(t^{j-i}\zeta)\prod_{u=1}^{i-j}\gamma(t^{-u}\zeta)^{-1}.
\]
The ratios $\gamma$ are found by the $q$-Petkov\v{s}ek step applied to
$\mathcal Q^*$, whose extreme coefficients are those of $\mathcal Q$, shifted.
(Equivalently, the elements correspond to the first-order left factors of
$\mathcal Q$, as the solutions correspond to its first-order right factors.)

\paragraph{The fit.} Given hypergeometric elements $w_i=L_ix$ and Laurent
polynomials $s_i,\rho\in K[\zeta,\zeta^{-1}]$, put
$\Pi=1-\sum_is_iL_i+\rho\,\mathcal Q=\sum_{j=0}^k\pi_jE^j$, with $E$ the shift.
Then
\[
  x_m=\sum_is_i(t^m)\,w_{i,m}+(\Pi x)_m\qquad(m\ge m_0).
\]
If every $\pi_j$ is regular at $\zeta=0$ and $\sum_j\pi_j(0)=0$, the remainder
$(\Pi x)_m$ tends to $0$ (Proposition~\ref{prop:remainder}), and the terms
$s_iw_i$ are channels whose sum has the same limit as $x$. These conditions are
linear in the Laurent coefficients of the $s_i$ and $\rho$: for a pole order
$e$, the coefficients of $\zeta^{-e},\dots,\zeta^{-1}$ in each $\pi_j$ vanish, and
the constant terms add up to $0$. The implementation lets $e$ run from $0$ to
$6$, and then drops elements one at a time while the conditions stay solvable,
so that few channels remain. Since a fit is only used to write $x$ as channels
plus a remainder, any solution is valid; the choice affects only whether the
channels then pass Proposition~\ref{prop:channels2}.

\paragraph{The test with remainders.} $S$ is additive over the diagonal blocks.
For each block, the section of its own $\beta$ is written as channels, exactly
(Proposition~\ref{prop:decomp}) or with a remainder. A fit found for $p=1$ is
restricted to the progression: $G(pm+r)$ has the ratio
$\prod_{l<p}g(t^{r+l}\zeta^p)$ in $m$. If the union of the channels of all
blocks, with similar classes merged, satisfies
Proposition~\ref{prop:channels2}, then $\beta_{pm+r}$ minus a sum of remainders
tends to $0$, so $S(\Delta)=0$. The implementation tries this for each period
after the plain channel test has failed for all combinations of readings.

\begin{example}\label{ex:sum-side}
The first list of the second set in Section~\ref{sec:automata} contains the
automata \texttt{[0R1 1U2 2R3 2U4 3R2 3U4 3U$\infty$ 4U3]} and
\texttt{[0R1 1U2 2R3 2U4 3R2 3U4 3U$\infty$ 4R3]}, numbers $7$ and $8$, with the
growth series $t\bigl((-t^2;t)_\infty-1\bigr)$. In reading \texttt{X}, the
$\beta$ of number~$7$ is a sum of channels along even $n$:
$\beta_{2m}=t/\bigl((1-t^2)(t^3;t^2)_{m-1}\bigr)-t$ for $m\ge1$. No reading of
number~$8$ has this property. In reading \texttt{Y} its $\beta$ satisfies, with $\zeta=t^n$,
\[
  (t^2+t^4\zeta)\beta_n-\bigl(t^2-(t^3-t^4)\zeta\bigr)\beta_{n+1}-(1+t^3\zeta)\beta_{n+2}+\beta_{n+3}=0,
\]
whose only $q$-hypergeometric solutions are $1$ and $(-t)^n$. The adjoint has the
solutions with ratios $1/t$ and $1/(t+t^4\zeta)$, which give the elements
\[
  w^{(1)}_n=-(t+t^{n+4})\beta_{n+1}+(t-1-t^{n+3})\beta_{n+2}+\beta_{n+3},\qquad
  w^{(2)}_n=-t\beta_{n+1}+(t-1)\beta_{n+2}+\beta_{n+3},
\]
with ratios $t$ and $t+t^{n+4}$. In both, the coefficients add up to $0$ at
$\zeta=0$, so both tend to $0$; the fit needs a pole, $s_1=-s_2=-1/((t^3+t^4)\zeta)$
and $\rho=0$, and then $\Pi=1-(tE+E^2)/(1+t)$. Thus
\[
  \frac{t\beta_{n+1}+\beta_{n+2}}{1+t}=s_1(t^n)\,w^{(1)}_n+s_2(t^n)\,w^{(2)}_n
  =t\bigl((-t^2;t)_{n+1}-1\bigr)\qquad(n\ge0),
\]
two channels with ratios $1$ and $1+t^{n+3}$, although $\beta_n$ itself is not a
sum of channels; the remainder $\beta_n-(t\beta_{n+1}+\beta_{n+2})/(1+t)$ has
valuation $2\lfloor n/2\rfloor+3$. Restricted to even $n$, the constants
cancel, and $t(-t^2;t)_{2m+1}$ pairs with the channel of number~$7$ through
Euler's witness $\chi=1-t^2\zeta^2$, $\lambda=2$, with a shifted form
($\Gamma=1+t\zeta^2+t^2\zeta^2+t^3\zeta^4$); this proves the difference.
\end{example}


\section{Euler's odd--distinct identity}\label{sec:euler}

\subsection{The input}
Let
\[
  h=1+Y(1-X-Y^2)^{-1}X,\qquad h'=1+Y(1-Y-XY)^{-1}X,\qquad\Delta=h-h'.
\]
The series $h-1$ is the sum of all words
$Y^{2i_1+1}XY^{2i_2}X\cdots Y^{2i_k}X$ with $k\ge1$. Moving each $X$ to the left
costs one factor $t$ per $Y$ it passes, so this word equals
$t^{|\lambda|}X^kY^{2i_1+\dots+2i_k+1}$ with
$\lambda=(2i_1+1,\,2i_1+2i_2+1,\,\dots,\,2i_1+\dots+2i_k+1)$, the partition into
odd parts read off from the number of $Y$'s preceding each $X$. Similarly the
words of $h'-1$, $Y\,w\,X$ with $w\in\{Y,XY\}^*$, encode the partitions into
distinct parts.

These expressions are not normalized realizations
\eqref{eq:normalization}: $1-X-Y^2\not\equiv1\pmod t$, and indeed
$(1-X-Y^2)^{-1}$ does not converge $t$-adically, since all the words $X^k$ have
weight $t^0$. Only the outer letters $Y\cdots X$ produce the weight. To move
them inside, let $\sigma$ be the automorphism $X\mapsto tX$, $Y\mapsto tY$ of
$A$. For a normal monomial $m=X^aY^b$, both $YmX$ and $tX\sigma(m)Y$ equal
$t^{a+b+1}X^{a+1}Y^{b+1}$; hence $YwX=tX\sigma(w)Y$ for every word $w$. Put
$P_{\mathrm o}=\sigma(1-X-Y^2)=1-tX-t^2Y^2$ and
$P_{\mathrm d}=\sigma(1-Y-XY)=1-tY-t^2XY$, both $\equiv1\pmod t$. As $w$
runs over $\{X,Y^2\}^*$, the words $\sigma(w)$ are the terms of the convergent
expansion of $P_{\mathrm o}^{-1}$, and likewise for $\{Y,XY\}^*$ and
$P_{\mathrm d}^{-1}$. Summing over words gives
\[
  h=1+tXP_{\mathrm o}^{-1}Y,\qquad h'=1+tXP_{\mathrm d}^{-1}Y.
\]
A normalized realization of $\Delta$ is
\[
  Q=\operatorname{diag}(P_{\mathrm o},P_{\mathrm d}),\qquad r=(X,X),\qquad e=(tY,-tY)^{\mathsf T},
\]
with compiled triple
\begin{equation}\label{eq:euler-triple}
  L=\begin{pmatrix}1-tT-t^2Y^2&0\\0&1-tY-t^2YT\end{pmatrix},\qquad
  f=\begin{pmatrix}1\\1\end{pmatrix},\qquad d=(tY,-tY).
\end{equation}
Writing $m_{\mathrm o}=\sum\eta_nY^n$ and $m_{\mathrm d}=\sum\zeta_nY^n$, coefficient
extraction gives $\eta_{2j+1}=0$ and
\[
  \eta_{2m}=\frac{t^{2m}}{(t;t^2)_{m+1}},\qquad
  \zeta_n=t^n(-t;t)_n,\qquad
  S(\Delta)=t\sum_{n\ge0}(\eta_n-\zeta_n).
\]
We will see that $\beta_{2m}=(t;t^2)_m^{-1}-(-t;t)_{2m}$, so
$S(\Delta)=(t;t^2)_\infty^{-1}-(-t;t)_\infty$: the identity $S(\Delta)=0$ is
Euler's theorem that partitions into odd parts and into distinct parts are
equinumerous \cite[26.10]{DLMF}.

\subsection{The identity: one inverse shift}
Equations \eqref{eq:reduced}, with the condition $(R\cdot f)(1)=0$, read
\[
  (1-tz)R_1(z)-t^2R_1(t^2z)=t,\qquad R_2(z)-t(1+tz)R_2(tz)=-t,\qquad R_1(1)+R_2(1)=0.
\]
(the last condition makes $\kappa=0$ in Proposition~\ref{prop:polycert}).
There is no polynomial solution: $-tzR_1(z)$
would be the unique top-degree term. The homogeneous equations
$(1-tz)h(z)=t^2h(t^2z)$ and $h(z)=t(1+tz)h(tz)$ have no rational solution, since
the divisor of $h(t^2z)/h(z)$, resp.\ $h(tz)/h(z)$, has degree $0$ on every orbit
$\alpha t^{\mathbb Z}$ while $1-tz$ and $1+tz$ have a single zero. So the rational
solution is unique, $R=(-T^{-1},T^{-1})$, with a pole of order $s=1$ at $T=0$, as
Proposition~\ref{prop:forced-pole} predicts. Left division by $1-Y$ gives
$P=((1+Y)T^{-1}-t,\,-T^{-1})$. In the shifted form \eqref{eq:shift},
\[
  \Lt=\begin{pmatrix}1-Y^2-tT&0\\0&1-Y-tYT\end{pmatrix},\quad \ft=f,\quad \dt=(tYT,-tYT),
\]
\begin{equation}\label{eq:euler-cert}
  \dt=(1-Y)P_s+R_s\Lt,\qquad P_s=(1+Y-tT,\,-1),\quad R_s=(-1,1),\quad R_s\ft=0,
\end{equation}
(e.g.\ $(1-Y)(1+Y-tT)-(1-Y^2-tT)=tYT$). Hence, with $\mu=T^{-1}m$,
\begin{equation}\label{eq:euler-b}
  b=(1+Y-tT)\mu_{\mathrm o}-\mu_{\mathrm d},\qquad (1-Y)b=C.
\end{equation}
The auxiliary series are integral but divergent: modulo $t$ they are
$(1-Y^2)^{-1}$ and $(1-Y)^{-1}$. Only the output $b$ will be shown to converge.

\subsection{Convergence: channels and a dilation witness}
Write $\mu_{\mathrm o}=\sum x_nY^n$, $\mu_{\mathrm d}=\sum y_nY^n$. From $\Lt\mu=f$,
\[
  (1-t^{n+1})x_n=x_{n-2}+\delta_{n0},\qquad y_0=1,\quad y_n=(1+t^n)y_{n-1}.
\]
The state $v_n=(x_n,x_{n-1},y_n)$ satisfies $v_{n+1}=A(t^n)v_n$, and
\eqref{eq:euler-b} gives $\beta_n=\pi(t^n)v_n$ with
\[
  A(z)=\begin{pmatrix}0&\frac1{1-t^2z}&0\\1&0&0\\0&0&1+tz\end{pmatrix},\qquad
  \pi(z)=(1-tz,\ 1,\ -1).
\]
The two-step product $A(tz)A(z)$ is diagonal. Along $n=2m$ the middle channel
vanishes ($x_{\mathrm{odd}}=0$), and the others are
\[
  U_m=(1-t^{2m+1})x_{2m}=\frac1{(t;t^2)_m},\qquad V_m=y_{2m}=(-t;t)_{2m},\qquad
  \beta_{2m}=U_m-V_m,
\]
with multipliers $g_U(\zeta)=(1-t\zeta^2)^{-1}$ and
$g_V(\zeta)=(1+t\zeta^2)(1+t^2\zeta^2)$, $U_0=V_0=1$, so the initial ratio is
$-1$. The ratio is
$\theta(\zeta)=(1-t\zeta^2)(1+t\zeta^2)(1+t^2\zeta^2)$; for $\lambda=2$ the
degree equation gives $\deg\chi=6/3=2$, and the recursion returns
\begin{equation}\label{eq:euler-witness}
  \chi(\zeta)=1-t^2\zeta^2,\qquad \chi(\zeta^2)\chi(t\zeta^2)=(1-t^2\zeta^4)(1-t^4\zeta^4)=\theta(\zeta)\chi(\zeta).
\end{equation}
Proposition~\ref{prop:channels}(b) applies with $i_1=c=2$:
\[
  \beta_{2m}=U_m\Bigl(1-\prod_{j=m}^{2m-1}(1-t^{2j+2})\Bigr),\qquad \val(\beta_{2m})\ge2m+2.
\]
The same channels, witness and bound are obtained from the scalar recurrence
of $\beta$ (Section~\ref{sec:intrinsic}), whichever realization of the two
series is used: PBW expressions, automata, or redundant automata.
By Corollary~\ref{cor:progression}, $S(\Delta)=0$, $b\in\Bhol\cap\Bhat$, and
$\val(\beta_n)\ge k_\Delta(n)=\lceil n/2\rceil$ for all $n$ (here $\alpha=1$, $D=2$).
Proposition~\ref{prop:canonical} then supplies the unique $a\in\Ahat$ with
$\Delta=(1-X)a+b(1-Y)$. The odd indices need no separate argument.

\begin{remark}[other routes]\label{rem:obstruction}
The row $\pi$ can also be found directly as a vector telescoper from the
unshifted system \eqref{eq:euler-triple}: the state
$v_n=(\eta_n,\zeta_n,\eta_{n-1})$ satisfies $v_{n+1}=A_0(t^n)v_n$ with
\[
  A_0(z)=\begin{pmatrix}0&0&\frac{t^2}{1-t^2z}\\0&t(1+tz)&0\\1&0&0\end{pmatrix},
  \qquad c_{n+1}=(t,-t,0)\,v_n,
\]
and the row $\pi_0(z)=\bigl((1-tz)/z,\,-1/z,\,t/z\bigr)$ satisfies
$\pi_0(tz)A_0(z)-\pi_0(z)=(t,-t,0)$ and $\pi_0(1)v_0=0$, so
$\beta_n=\pi_0(t^n)v_n$ without any operator certificate. Uniform contraction cannot
prove its decay: for the initial vector $((1-t)^{-1},0,0)$ the same
recurrence and row yield the unit sequence $U_m$ at even indices, so no bound valid for all initial vectors exists.
The decay is a property of the selected initial values, which is exactly what
condition~(i) of Proposition~\ref{prop:channels} records.
\end{remark}

\section{Readings: other series with the same sum}\label{sec:readings}

The test of Section~\ref{sec:intrinsic} depends only on the partial sums $\beta$
of the marginal $C$. The marginal, however, depends on the realization, while
$S(\Delta)=C(1)$ does not. Partition generating functions typically have two
marginals of different nature. For partitions into distinct parts $\ge2$, with
sum $S=(-t^2;t)_\infty-1$,
\[
  C_{\mathrm{lp}}=\sum_{n\ge2}t^n(-t^2;t)_{n-2}\,Y^n,\qquad
  C_{\mathrm{np}}=\sum_{k\ge1}\frac{t^{k(k+3)/2}}{(t;t)_k}\,Y^k
\]
count them by largest part and by number of parts. The partial sums of
$C_{\mathrm{lp}}$ are products, $\beta_n=(-t^2;t)_{n-1}-1$, with
$\val(S-\beta_n)=n+1$. Those of $C_{\mathrm{np}}$ are partial sums of a
$q$-hypergeometric series, with $\val(S-\beta_k)=(k+1)(k+4)/2$. When the two
sides of $\Delta$ are realized in these two ways, $C\not\equiv0$ and the channel
test finds no proof (systems~1 and~2 of Section~\ref{sec:six} are an example),
although $S(\Delta)=0$. This section replaces each block of the marginal system
by other systems with the same sum.

\subsection{The second marginal}
Let $\rho_Y$ be evaluation at $Y=1$, followed by renaming $X$ to $Y$, and let
$\Psi:A\to\Op$ be the isomorphism $\Psi(X^aY^b)=Y^aT^b$. On normal monomials of
$\Op$ put $\operatorname{sw}(Y^bT^a)=Y^aT^b$, extended linearly.

\begin{lemma}[transposition]\label{lem:transpose}
Let $(L,f,d)$ be a normalized marginal system with sum $S$. Then
\[
  L'_{ij}=\operatorname{sw}(L_{ji}),\qquad
  f'_i=\operatorname{sw}(d_i)\big|_{T=1},\qquad
  d'_j=f_j(T)
\]
is a normalized marginal system with the same sum $S$.
\end{lemma}
\begin{proof}
Put $Q=\Phi^{-1}(L^{\mathsf T})$, $e=\Phi^{-1}(d)$ and $r_j=f_j(Y)\in A$. By
\eqref{eq:triple}, $(L,f,d)$ is the marginal system of the right realization
$wQ=r$, $\Delta=we$, so $S(\Delta)=S$. The same element is $\Delta=rw'$ with
$Qw'=e$, a \emph{left} realization, and $w'=Q^{-1}e$ is again given by
geometric inversion. For $q=X^aY^b$ and $w=X^iY^j$ we have
$qw=t^{bi}X^{a+i}Y^{b+j}$, hence $\rho_Y(qw)=\Psi(q)\cdot\rho_Y(w)$. Applying
$\rho_Y$ to $Qw'=e$ and to $\Delta=rw'$ gives $L'\cdot m'=f'$ and
$\rho_Y(\Delta)=d'\cdot m'$, with $m'=\rho_Y(w')$, $L'=\Psi(Q)$,
$f'=\rho_Y(e)$ and $d'=\Psi(r)$; these are the formulas above. Both marginals of
$\Delta$ take the value $S(\Delta)$ at~$1$.
\end{proof}

In the encoding of Section~\ref{sec:euler} a partition $\lambda$ becomes
$t^{|\lambda|}X^kY^j$ with $k$ parts and largest part $j$. So the two marginals
of such a realization count by largest part and by number of parts.

\subsection{Canonical realizations}\label{sec:canonical-real}
The second marginal depends on the system, not only on $C$. A system determined
by $C$ comes from an annihilator. Let $H\cdot C=g\in K[Y]$ as in
Section~\ref{sec:intrinsic}, write $H=\sum_bh_b(T)\,Y^b$ with coefficients on the
left, and let $h\in K[T]$ be the greatest common divisor of the $h_b$, the
\emph{content}: $H=h(T)H'$.

\begin{proposition}[canonical realization]\label{prop:canonical-real}
The series $u=H'\cdot C$ is a polynomial: $h(t^n)u_n=g_n$, so $u_n=0$ for
$n>\deg g$ except at the finitely many $n$ with $h(t^n)=0$, and $u$ is computed
from finitely many $c_n$. If $t^{-v}H'\equiv c\pmod t$ for some $v\in\mathbb Z$
and $c\in\Q^\times$, then
\[
  \bigl(t^{-v}H'/c,\ \ t^{-v}u/c,\ \ 1\bigr)
\]
is a normalized $1\times1$ marginal system whose marginal is $C$. If $H$ has
least $Y$-degree among all $H$ with $H\cdot C\in K[Y]$, the transpose of this
system (Lemma~\ref{lem:transpose}) depends only on $C$.
\end{proposition}
\begin{proof}
$h(T)$ acts on $Y^n$ as multiplication by $h(t^n)$, and the zeros of $h$ at
$z=t^n$ are bounded as in Section~\ref{sec:convergence}. The system is
normalized and $C\in\Bhat$ solves it, so $C$ is its selected solution. The $H$
with $H\cdot C\in K[Y]$ form a left ideal of $\Op$. If $H_1,H_2$ are elements of
least degree, a combination $b(T)H_1-a(T)H_2$ of lower degree lies in the ideal,
so it vanishes: $H_1$ and $H_2$ are proportional over $K(T)$, and their
content-free parts differ by a factor $\lambda\in K^\times$. Replacing $(H',u)$
by $(\lambda H',\lambda u)$ multiplies $L'$ and $d'$ of the transpose by
$\lambda$ and leaves its marginal $d'\cdot(L')^{-1}\cdot1$ unchanged.
\end{proof}

The implementation takes $H$ from \texttt{annihilator}, of least degree under
a common bound on the degrees of $H$ and $R$. Only the last statement of the
proposition depends on minimality; every reading has the sum $S$ in any case.

\subsection{Searching the readings}\label{sec:search}
If $\Delta=h-h'$ comes from two realizations, its marginal system is block
diagonal; in general the blocks are the connected components of the nonzero
pattern of $L$. The sum is additive over blocks, so every block can be replaced
by any of its \emph{readings}: the block itself (\texttt{Y}), its transpose
(\texttt{X}), and the transposes of the canonical realizations of these two
(\texttt{cY$'$}, \texttt{cX$'$}), when they exist. Every combination is a
normalized system with the same sum, and all previous methods apply to it
unchanged: a violated a-priori rate proves $S(\Delta)\ne0$, and a zero
certificate, a rational certificate without pole at $T=0$, a polynomial
certificate or a channel proof proves $S(\Delta)=0$.
The annihilator of a combination is assembled from those of its blocks: if
$H_k\cdot C_k=g_k$ and $H=A_1H_1=A_2H_2=\cdots$ is a least common left multiple,
then $H\cdot\sum_kC_k=\sum_kA_kg_k$. Readings of one series may also differ by a
power of $Y$: the second marginal of a realization counts the parts from $0$ or
from $1$, depending on where the outer letters are placed. Since
$(Y^sC_k)(1)=C_k(1)$, when the marginals of the blocks start at different powers
of $Y$ and cancel on their first coefficients after multiplying each block by
the power of $Y$ that aligns them ($d\mapsto Y^sd$), the aligned combination is
tried as well, and the zero test decides it. The implementation runs the cheap
tests on the system as given, then on all combinations (at most four readings
per block), and then the channel test with increasing period across all
combinations.

\subsection{Six realizations of one series}\label{sec:six}
The following six systems all have the sum
$S=(-t^2;t)_\infty-1=t^2+t^3+t^4+2t^5+2t^6+3t^7+\cdots$ of partitions into
distinct parts $\ge2$; in each, $f=e_1$.
\[
\begin{array}{@{}c@{\ \ }l@{\qquad}c@{\ \ }l@{}}
1&L=\begin{pmatrix}1-tT&-tT\\-tY&1\end{pmatrix},\ d=(0,tY)&
2&L=\begin{pmatrix}1&-tT\\-tY&1-tY\end{pmatrix},\ d=(0,tY)\\[12pt]
3&L=\begin{pmatrix}1&-tY\\-tT&1-tT\end{pmatrix},\ d=(0,tY)&
4&L=\begin{pmatrix}1&-tY&0\\-tY&1&-tY\\-tT&-tT&1\end{pmatrix},\ d=(0,tY,0)\\[18pt]
5&L=\begin{pmatrix}1&-tY&0\\-tY&1&-tT\\-tT&-tT&1\end{pmatrix},\ d=(0,tY,0)&
6&L=\begin{pmatrix}1&0&0\\0&1-tY&-tY\\-tT&-tT&1\end{pmatrix},\ d=(0,tY,tY)
\end{array}
\]
Their readings are of four kinds, distinguished by the valuations $\val(c_n)$:
by largest part (\textsf{lp}, $\val(c_n)$ linear in $n$), by number of parts
(\textsf{np}, quadratic), with alternating parity (\textsf{alt}:
$2,5,4,7,6,9,\dots$), and supported on even $n$ (\textsf{ev}):
\begin{center}
\begin{tabular}{@{}lcccccc@{}}
\toprule
reading & 1 & 2 & 3 & 4 & 5 & 6\\\midrule
\texttt{Y}    & \textsf{np} & \textsf{lp} & \textsf{np} & \textsf{alt} & \textsf{alt} & \textsf{lp}\\
\texttt{X}    & \textsf{lp} & \textsf{np} & \textsf{lp} & quadratic & \textsf{ev} & \textsf{np}\\
\texttt{cY$'$} & -- & \textsf{np} & -- & \textsf{np} & -- & \textsf{np}\\
\texttt{cX$'$} & \textsf{np} & \textsf{lp} & \textsf{np} & \textsf{lp} & \textsf{lp} & --\\
\bottomrule
\end{tabular}
\end{center}
(--: no canonical realization, since $t^{-v}H'\not\equiv c\pmod t$.) All
fifteen differences $\Delta_i-\Delta_j$ are proved, with the readings of
$\Delta_i$ and $\Delta_j$ shown: by a rational certificate without pole at $T=0$
for the least annihilator (R), by the zero test (Z: the two readings are equal
series), or by the channel test along $n=2m+r$ (C$_r$):
\begin{center}
\small
\begin{tabular}{@{}c|ccccc@{}}
 & 2 & 3 & 4 & 5 & 6\\\hline
1 & R\ \texttt{X},\texttt{Y} & R\ \texttt{Y},\texttt{Y} & R\ \texttt{X},\texttt{Y} & C$_0$\ \texttt{X},\texttt{X} & R\ \texttt{X},\texttt{Y}\\
2 & & R\ \texttt{X},\texttt{Y} & R\ \texttt{Y},\texttt{Y} & C$_0$\ \texttt{Y},\texttt{X} & R\ \texttt{Y},\texttt{Y}\\
3 & & & R\ \texttt{X},\texttt{Y} & C$_0$\ \texttt{X},\texttt{X} & Z\ \texttt{X},\texttt{Y}\\
4 & & & & C$_0$\ \texttt{cX$'$},\texttt{X} & R\ \texttt{Y},\texttt{Y}\\
5 & & & & & C$_0$\ \texttt{X},\texttt{Y}
\end{tabular}
\end{center}
In the nine R cases the two readings are different series, and their
difference has a least annihilator of order $1$ or $2$, against $2$ or $3$ for
the annihilator of the system; for 1--3, for instance,
$R=(t^2+t)/\bigl((t+1)^2T^2-(t^2+2t)T+t\bigr)$.

In the five channel proofs, $\beta_{2m}$ is a combination of two product
channels with ratios $1/(1-t^3\zeta^2)$ and $(1+t^{a}\zeta^2)(1+t^{a+1}\zeta^2)$,
whose quotient is $\theta(\zeta)=(1-t^3\zeta^2)(1+t^a\zeta^2)(1+t^{a+1}\zeta^2)$.
For $a\in\{2,3\}$ ($a=3$ for 1--5 and 4--5, $a=2$ for 3--5 and 5--6) the witness
is $\chi=1-t^{2a-2}\zeta^2$ with $\lambda=2$ and offset $\mu=1$:
\[
  \chi(t\zeta^2)\,\chi(t^2\zeta^2)=(1-t^{2a}\zeta^4)(1-t^{2a+2}\zeta^4)=\theta(\zeta)\,\chi(\zeta).
\]
For $a=1$ (pair 2--5) there is no witness for $\theta$ itself, and the proof
uses the shifted form of Section~\ref{sec:intrinsic}: $\Gamma=1-t\zeta^2$ and
$\chi=1-t^2\zeta^2$ with $\mu=0$. With the systems as given (the readings
\texttt{Y} alone), only the four differences 1--3, 2--4, 2--6 and 4--6 are
proved.

\subsection{Four families}\label{sec:families}
Three more families were treated in the same way. Together with the six systems
above (family A) there are $27$ systems in four families, with one sum per family:
\[
  S_A=(-t^2;t)_\infty-1,\qquad S_B=(-t;t)_\infty-1-t,\qquad
  S_C=t\,(-t;t)_\infty,\qquad S_D=t\,\bigl((-t;t)_\infty-1\bigr).
\]
In all of them $f=e_1$; the matrices $L$ are written row by row, as in the input
of the implementation:
\begin{center}
\footnotesize
\begin{tabular}{@{}l@{\ \ }l@{\ \ }l@{\quad}l@{\ \ }l@{\ \ }l@{}}
B1 & $[1,0,0;\;-tT,1-tT,-tY;\;-tY,-tY,1]$ & $(0,tY,0)$ & B2 & $[1,0,0;\;-tY,1-tY,-tY;\;-tT,-tT,1]$ & $(0,tY,0)$ \\
B3 & $[1,0,0;\;-tY,1-tY,-tT;\;-tT,-tT,1]$ & $(0,tY,0)$ & B4 & $[1-tT,-tT,0;\;-tY,1,0;\;0,-tY,1]$ & $(0,tY,tY)$ \\
B5 & $[1,0,0;\;-tT,1-tY,-tY;\;0,-tT,1]$ & $(0,tY,tY)$ & B6 & $[1,0,0;\;0,1,-tY;\;-tT,-tT,1-tT]$ & $(0,tY,tY)$ \\
\addlinespace
C1 & $[1-tY,-tY,0;\;0,1,-tY;\;-tT,-tT,1]$ & $(tY,0,0)$ & C2 & $[1-tY,0,-tY;\;-tT,1,-tT;\;0,-tT,1]$ & $(tY,0,0)$ \\
C3 & $[1,0,0;\;-tT,1-tT,-tY;\;0,-tY,1]$ & $(tY,tY,0)$ & C4 & $[1,0,0;\;-tT,1-tY,-tY;\;0,-tT,1]$ & $(tY,tY,0)$ \\
C5 & $[1,0,0;\;-tT,1-tY,-tT;\;0,-tT,1]$ & $(tY,tY,0)$ & C6 & $[1,0,-tY;\;-tY,1,0;\;-tT,0,1-tT]$ & $(tY,tY,0)$ \\
C7 & $[1,0,0;\;-tY,1,-tY;\;-tT,-tY,1-tT]$ & $(tY,tY,0)$ & C8 & $[1,0,0;\;-tT,1,-tY;\;0,-tT,1-tT]$ & $(tY,tY,tY)$ \\
\addlinespace
D1 & $[1-tT,-tY;\;-tY,1]$ & $(0,tY)$ & D2 & $[1,-tT;\;-tT,1-tY]$ & $(0,tY)$ \\
D3 & $[1,0,0;\;-tT,1-tT,-tY;\;0,-tY,1]$ & $(0,tY,0)$ & D4 & $[1,0,0;\;-tT,1-tY,-tY;\;0,-tT,1]$ & $(0,tY,0)$ \\
D5 & $[1,0,-tY;\;-tY,1-tY,0;\;-tT,-tT,1]$ & $(0,tY,0)$ & D6 & $[1,0,-tT;\;-tY,1-tY,0;\;-tT,-tT,1]$ & $(0,tY,0)$ \\
D7 & $[1,0,0;\;-tT,1,-tY;\;0,-tT,1-tT]$ & $(0,tY,tY)$ & && \\
\end{tabular}
\end{center}
All $79$ differences within a family are proved:
\begin{center}
\begin{tabular}{@{}llccccc@{}}
\toprule
family & sum & systems & differences & zero test & rational & channels\\\midrule
A & $(-t^2;t)_\infty-1$ & 6 & 15 & 1 & 9 & 5\\
B & $(-t;t)_\infty-1-t$ & 6 & 15 & 1 & 6 & 8\\
C & $t\,(-t;t)_\infty$ & 8 & 28 & 0 & 12 & 16\\
D & $t\,((-t;t)_\infty-1)$ & 7 & 21 & 2 & 7 & 12\\\midrule
all & & 27 & 79 & 4 & 34 & 41\\
\bottomrule
\end{tabular}
\end{center}
The $41$ channel proofs are product identities in the sense of
Proposition~\ref{prop:forced-pole}. In the readings used, no annihilator of
lower order than that of the system (order $3$) is found, the least pole at
$T=0$ of a rational certificate has order~$2$, and $\beta_{2m}$ consists of two
product channels that cancel through one dilation witness, $\chi=1-t^2\zeta^2$
($38$ pairs) or $1-t^4\zeta^2$ ($3$), with $\lambda=2$ and offset $\mu=0$ ($33$)
or $\mu=1$ ($8$); $20$ of them need a shifted form ($\Gamma\ne1$).
In the other $38$, the readings found first, all of them \texttt{X} or
\texttt{Y}, are equal series (zero test, $4$ pairs), differ by a nonzero
polynomial ($2$ pairs, $A_C=1$), or differ by a series whose least annihilator
has order~$1$ ($23$ pairs) or~$2$ ($9$ pairs); the annihilator of the system has
order up to~$4$. Every rational certificate found is without pole at $T=0$. For
B1--B3, for instance, the readings (\texttt{X},\,\texttt{Y}) differ by a
$q$-hypergeometric series, the system gives an annihilator of order~$2$, and
the certificate for $A_C$ is
\[
  R=\frac{t^3}{(t^2T^2+t-1)(T^2+t-1)},\qquad R(0)=\frac{t^3}{(t-1)^2}.
\]
A canonical reading occurs in one proof (A4--A5). With the systems as given,
without other readings, $23$ of the $79$ differences are proved.

\subsection{Lists of automata with a common growth series}\label{sec:automata}
A larger test consists of $497$ lists of automata over two letters, with
conjecturally the same growth series within each list. An automaton is given by
its transitions $iVj$, from state $i$ to state $j$ reading $V\in\{R,U\}$, with
initial state $0$ and final state $\infty$; for instance
\texttt{[0R1 1R2 1U1 1U$\infty$ 2R1]} recognizes $R(U+R^2)^*U$. With $R=Y$ and $U=X$,
an automaton stands for the sum of its accepted words, and its growth series is
$S$ of that sum: the generating function of the words by the number of pairs
$R\cdots U$ (Euler's odd series in the example). Every accepted word has the form
$R\cdots U$, so the sum is $Y(\sum_ww)X$ over the inner words, realized as in
Section~\ref{sec:euler} (\texttt{automaton} reads this notation). The lists
contain $1116$ automata with two or three inner states, in $422$ pairs and $75$
lists of three to eight; within each list the series agree up to $t^{40}$. That
makes $813$ differences.

All $813$ differences are proved directly (\texttt{automata\_quantumwz.jl}), with
a time limit of $300$ seconds for the channel search:
\begin{center}
\begin{tabular}{@{}lr@{}}
\toprule
rational certificate without pole at $T=0$ & $728$\\
zero test (two after aligning the readings by a factor $Y$) & $14$\\
channel test, period $2$, $\lambda=2$ & $69$\\
channel test, period $3$, $\lambda=3$ & $2$\\\midrule
all & $813$\\
\bottomrule
\end{tabular}
\end{center}
Most proofs use the readings (\texttt{X},\,\texttt{Y}) ($485$) or
(\texttt{Y},\,\texttt{Y}) ($296$); canonical readings occur in five channel
proofs. Of the $71$ channel proofs, $39$ need a shifted form $\theta_w$
($\Gamma\ne1$), as for the boundary factor in Euler's identity of
Section~\ref{sec:intrinsic}; the witnesses are $1-t^2\zeta^2$ ($66$),
$1-t^4\zeta^2$ ($3$) and $1-t^3\zeta^3$ ($2$, with period $3$). The median time per
difference is $0.2$ seconds; $11$ take more than ten seconds, the longest
two and a half minutes, and all $813$ take $17$ minutes. Three changes to the method
came from this data: the least annihilator is reconstructed from modular images
(the direct solution over $K$ took minutes for systems of $20$ unknowns),
stationary channels may differ by a boundary factor (the shifted forms
$\theta_w$), and readings may differ by a power of $Y$ (the aligned
combinations).

\paragraph{A second set.} A second file has $49$ lists of $10$ to $19$ automata
with two to four inner states: $579$ automata and $3237$ differences. Run with
the remainder test of Section~\ref{sec:remainders} and a time limit of $150$
seconds, all but $29$ differences are proved directly:
\begin{center}
\begin{tabular}{@{}lr@{}}
\toprule
rational certificate without pole at $T=0$ & $2726$\\
zero test & $70$\\
channel test, period $2$, $\lambda=2$ & $399$\\
channel test with a remainder, period $2$, $\lambda=2$ & $13$\\
not proved (Gauss's identity, below) & $29$\\\midrule
all & $3237$\\
\bottomrule
\end{tabular}
\end{center}
The remainder test was added for this set. In the first list, automaton~$8$ is
the sum side of Example~\ref{ex:sum-side}; its differences with numbers $7$ and
$10$ were not proved before, and the list was certified only through the other
differences. All $13$ differences of this kind, in four lists, are now proved,
each with a fit of pole order~$1$ for one block, found for $p=1$ and restricted
to a progression of period~$2$. Of the $412$ channel proofs, $273$ need a shifted form; the
witnesses are $1-t^2\zeta^2$ ($401$) and $1-t^4\zeta^2$ ($11$). Most proofs use
the readings (\texttt{Y},\,\texttt{Y}) ($2485$) or (\texttt{X},\,\texttt{Y})
($548$); canonical readings occur in $10$. The median time of a proof is $0.1$
seconds and the longest $28$ seconds; the proved differences take $44$ minutes
in all, and the $29$ others $94$ minutes, until the time limit (one step of the
search can overrun it by a few minutes). So $47$ of the $49$ lists are
certified.

\paragraph{Gauss's identity.} The $29$ unproved differences are those of
automaton~$10$ of list~$34$, and of automata $6$ and $7$ of list~$42$, with the
other members of their lists (the difference of $6$ and $7$ is proved, by the
zero test). Automaton~$10$ of list~$34$ is
\texttt{[0R1 1R2 1U3 1U$\infty$ 2R1 2U2 3R4 3U4 3U$\infty$ 4R3]}. Its words are
$R\,w_1\cdots w_k\,U$ and $R\,w_1\cdots w_k\,U\,v_1\cdots v_l\,U$ with factors
$w_i=RU^{a_i}R$ and $v_j\in\{UR,RR\}$. The letters $U$ of $w_i$ follow $2i$
letters $R$, so the factors $w_i$ contribute $1/(t^2;t^2)_k$, and $N=2k+1$
letters $R$ precede the next $U$. An ending with $u$ factors $UR$ and $r$ factors
$RR$ has the total weight $t^{N(u+2)+u(u+1)/2+2r}\binom{u+r}{u}_{t^2}$. Summing
over $r$, then over $k$ by $\sum_kz^k/(t^2;t^2)_k=1/(z;t^2)_\infty$,
\[
  S=\frac{t}{(t^2;t^2)_\infty}\Bigl(1+\sum_{u\ge0}t^{\binom{u+2}2}\Bigr)
   =\frac{t\,\psi(t)}{(t^2;t^2)_\infty},\qquad \psi(t)=\sum_{v\ge0}t^{\binom{v+1}2},
\]
while the other members of the list have $S=t(-t;t)_\infty$. The same
computation gives $S=\psi(t)/(t^2;t^2)_\infty-1$ for automaton~$6$ of list~$42$,
\texttt{[0R1 1R2 1U3 1U$\infty$ 2R1 2U2 2U$\infty$ 3R4 4R3 4U3 4U$\infty$]}, whose
list has $S=(-t;t)_\infty-1$. These equalities are Gauss's identity
$\psi(t)=(t^2;t^2)_\infty/(t;t^2)_\infty$, a theta-function identity. None of
the three automata has a reading in which $\beta$ is a sum of channels, exactly or
up to a remainder, for $p\le2$: in reading \texttt{X}, the partial sums of
automaton~$10$ have the $q$-hypergeometric constituents $1$ and $1/(t^4;t^2)_n$
and a third part that is not a sum of channels, and the hypergeometric elements
of their module give no fit (pole order up to $10$). So lists $34$ and $42$ are
the two lists that the tests do not certify. This is a limitation of the
certificate classes: they prove identities between products, and sum sides that
a finite identity within the same module relates to a product (as in
Example~\ref{ex:sum-side}), but not a theta series such as~$\psi$.

\section{Implementation}\label{sec:implementation}
\sloppy

The accompanying module \texttt{QuantumWZ.jl} implements the route of
Sections~\ref{sec:intrinsic}, \ref{sec:remainders} and~\ref{sec:readings}. Both algebras are
\textsc{Oscar} PBW algebras \cite{OscarPBW} over $K=\Q(t)$: \texttt{A} with
relation $YX=tXY$ and \texttt{O} with $TY=tYT$, with normal monomials $X^aY^b$
and $Y^bT^a$ as in the paper. Operators stay PBW elements throughout and are read
by a single function, \texttt{pbw\_terms}. A power series is stored as its
vector of coefficients, and the action of $\Op$ is used only through the formula
$[Y^k](cY^bT^a\cdot v)=c\,t^{a(k-b)}v_{k-b}$ (\texttt{ycoeff}). Recurrences live
in $F=\Q(t,z)$, the fraction field of the Flint polynomial ring $\Q[t,z]$, where
$z$ stands for $t^n$ (or for $T$); factorization over $K$ is factorization in
$\Q[t,z]$. Linear algebra is plain Gauss--Jordan elimination on matrices of
normalized fractions (\texttt{rref\_field}, \texttt{kernel\_field},
\texttt{solve\_field}), over $K$ whenever the entries are free of $z$; on the
annihilator systems it is faster than the generic fraction-free routines by two
to three orders of magnitude.

\begin{enumerate}
\item \texttt{compile} produces $(L,f,d)$ by \eqref{eq:triple}, and
  \texttt{difference} the block-diagonal system of $\Delta_1-\Delta_2$ from two
  realizations or two compiled systems; \texttt{automaton} also reads the
  notation of Section~\ref{sec:automata}. \texttt{attempt} splits the system
  into \texttt{blocks} and runs the following steps (\texttt{cheap\_stages}) on
  the system as given; if nothing is decided, it lists the \texttt{readings} of
  each block (\texttt{transpose\_system}, \texttt{canonical}) and runs them on
  every other combination, together with the aligned combinations of
  Section~\ref{sec:search} (\texttt{aligned}).
\item \texttt{solve\_series} and \texttt{marginal\_coeffs} give $c_0,\dots,c_N$
  and $\beta_0,\dots,\beta_N$; $N$ is extended automatically when a proof needs
  more terms. \texttt{nonzero\_witness} applies Proposition~\ref{prop:rate}(1).
\item \texttt{annihilator} computes $H$ and $R$ with $Hd=RL$ over $K(T)$, of
  least degree under a common degree bound, and checks the identity in
  \texttt{O}. A block-diagonal system is done block by block, with $H$ a least
  common left multiple of the block annihilators (\texttt{lclm}). If
  $g=\eps(Rf)=0$, \texttt{exceptional\_set} and Proposition~\ref{prop:zero}
  decide $C\equiv0$ (\texttt{:zero\_marginal}).
\item \texttt{reduce\_annihilator} looks for a right factor $A$ of $H$ of lower
  order with $A\cdot C=g_A$ a polynomial (Proposition~\ref{prop:factor}). The
  coefficients of $A$, polynomials in $T$, are guessed from the $c_n$: the least
  order and then the least degree are found on modular images ($t\mapsto T_0$ in
  $\mathrm{GF}(2^{61}-1)$, \texttt{marginal\_mod}) with twice as many equations
  as unknowns. The kernel, one-dimensional, has entries in $\Q(t)$ but is
  computed from its images at further points $t\mapsto\tau$: each entry is
  interpolated and rationally reconstructed in $t$, then its coefficients are
  lifted to $\Q$ (\texttt{reconstruct\_kernel}); only for small systems is it
  solved over $K$ directly, since the $c_n$ have large $t$-degree. The guess
  is then proved by \texttt{check\_factor}: right division $H=QA$, the
  identity $\widetilde QA=DH$ in \texttt{O}, $\widetilde Q\cdot g_A=D\cdot g$,
  and the exceptional indices of $\widetilde Q$. The search is limited to $32$
  unknowns; beyond it, or if the reconstruction fails, $A=H$.
\item \texttt{rational\_certificate} solves \eqref{eq:scalar-cert} with $A$ in
  place of $H$ for rational $R$: \texttt{universal\_denominator}, the pole bound at $0$ and the
  degree bound, then the numerator by forward substitution in its coefficients
  (\texttt{polynomial\_solution}: the coefficient of $z^{v+i}$ determines the
  $i$-th unknown except where $\chi(t^i)=0$, so only those few free parameters
  are found by linear algebra). Without a pole at $T=0$ it returns
  \texttt{:rational}, or \texttt{:nonzero} with the exact value
  $S(\Delta)=\kappa\ne0$. Inside \texttt{attempt} the search is skipped when the
  cleared denominator has degree above~$120$. \texttt{polynomial\_certificate}
  looks for a polynomial certificate of the system itself (\texttt{:polynomial}).
\item For $p\le p_{\max}$, every combination and every residue,
  \texttt{progression\_test}
  \begin{itemize}
  \item turns $A(1-Y)$ into \eqref{eq:beta-rec} (\texttt{beta\_recurrence}) and
    its section \eqref{eq:section-rec} (\texttt{section\_recurrence});
  \item finds the $q$-hypergeometric solutions (\texttt{qhyper}); a candidate is
    rejected from the extreme coefficients of the $T_i$ alone, before any
    polynomial is formed;
  \item fixes the coefficients $a_i$ from $k$ values (\texttt{decompose},
    Proposition~\ref{prop:decomp}) and merges similar channels
    (\texttt{merge\_similar}, \texttt{similarity});
  \item applies Proposition~\ref{prop:channels2} (\texttt{channel\_proof}), with
    dilation witnesses for every admissible $\lambda$ and offsets $|\mu|\le6$
    (\texttt{dilation\_witness}), for $\theta$ itself and for its shifted forms
    $\theta_w$ (\texttt{shift\_forms}).
  \end{itemize}
  Success gives \texttt{:channels}, with the progression, the channels, the
  witnesses and the bound $\val(\beta_{pm+r})\ge sm+\kappa$.
\item For each period, once \texttt{progression\_test} has failed for every
  combination, \texttt{remainder\_test} runs the test with remainders of
  Section~\ref{sec:remainders}. Each block gets its own least annihilator and
  $\beta$ recurrence (\texttt{block\_reading}); \texttt{block\_channels} writes
  the section of its $\beta$ as channels, exactly or with a remainder
  (\texttt{hyper\_elements}: \texttt{qhyper} on the adjoint and the
  concomitant; \texttt{remainder\_fit}: the Laurent conditions;
  \texttt{restrict}: a fit for $p=1$ along the progression). The channels of
  all blocks are merged (\texttt{merge\_channels}) and passed to
  \texttt{channel\_proof}. Block data are cached across combinations. Success
  gives \texttt{:channels} with the field \texttt{remainders}: for each block,
  \texttt{:none} or the pole order of its fit.
\end{enumerate}
Every result records the readings used. The result \texttt{:unknown} is not a
disproof.

\paragraph{Tests.} \texttt{test\_quantumwz.jl} checks:
\begin{itemize}
\item $X-Y$ (\texttt{:rational}, no denominator), the second system of
  Example~\ref{ex:incomplete} (\texttt{:rational}, denominator $1+T^2$), and
  $\Delta=1$ (\texttt{:nonzero} at $n=0$);
\item Example~\ref{ex:incomplete} and Example~\ref{ex:two-realizations}
  (\texttt{:zero\_marginal}), together with further pairs of equal series given
  by different realizations (PBW, automata, a redundant automaton);
\item Euler's identity for four pairs of realizations (PBW and PBW, automaton
  and automaton, PBW and automaton, automaton and redundant automaton) and for
  $-\Delta$. All give the proof of Section~\ref{sec:euler}: readings
  (\texttt{Y},\,\texttt{Y}), period~$2$, a single dilation pair with $\lambda=2$
  and $\chi=1-t^2z^2$, and $\val(\beta_{2m})\ge2m+2$;
\item all fifteen differences of the six systems of Section~\ref{sec:six}, the
  readings and witness of the pair 4--5, the false identity ``distinct parts
  $\ge2$ versus distinct parts $\ge1$'' (\texttt{:nonzero}), and one difference
  entered with integer entries, as in $L=[-t*T+1\ \ -t*T;\ -t*Y\ \ 1]$;
\item the least annihilator: B1--B3 of Section~\ref{sec:families}
  (\texttt{:rational} without pole, with $A$ of order~$1$ where the system gives
  order~$2$), and C3--C5, whose readings differ by a polynomial ($A=1$);
\item automata in the notation of Section~\ref{sec:automata}: Euler's identity,
  a product identity that needs a shifted form ($\Gamma\ne1$), and a pair
  proved after aligning the readings by a factor~$Y$;
\item the sum side of Example~\ref{ex:sum-side}: the two hypergeometric elements,
  the fit with a pole of order~$1$, and the difference of numbers $7$ and $8$,
  proved with a remainder for number~$8$ only;
\item Glaisher's identity for $k=3$ (parts not divisible by $3$ versus parts
  used at most twice), with both sides given by automata: period~$3$ and a
  dilation witness with $\lambda=3$; and a false identity (parts not divisible by
  $3$ versus distinct parts), which ends with \texttt{:nonzero}.
\end{itemize}
\texttt{families\_quantumwz.jl} runs all $79$ differences of the four families of
Section~\ref{sec:families}, and \texttt{automata\_quantumwz.jl} all differences
within each list of a file in the notation of Section~\ref{sec:automata}; it
also reports whether the proved differences connect every list. \texttt{stress\_quantumwz.jl} draws random automata over
$\{X,Y\}$ and runs \texttt{attempt} on differences whose answer is known: two
unrelated automata (a proof would be wrong whenever $S\ne0$ is certain from
$\beta_N$), and an automaton against its transpose, against itself with a state
split in two, and against the transpose of its canonical realization ($S=0$ by
construction; a refutation would be wrong).

\paragraph{Limitations.} Channels must have ratios in $K$ (roots of the
coefficient equations are sought in $\Q(t)$), the degree of $C$ in
\texttt{qhyper} is capped at~$64$, periods are tried up to $p_{\max}$, offsets
up to $|\mu|\le6$, shifted forms over four windows, and stationary channels
must cancel in pairs through a dilation witness. The enumeration of candidates
in \texttt{qhyper} is
exponential in the number of irreducible factors of $q_0$ and $q_k$; above
$50\,000$ candidates it is skipped. A failing search tries every combination of
readings for every period, so \texttt{:unknown} can take long; the keyword
\texttt{time\_limit} of \texttt{attempt} bounds the channel search (it is
checked between progressions, and inside \texttt{qhyper} and
\texttt{section\_recurrence}). A fit has pole order at most~$6$, and its
elements come only from $q$-hypergeometric solutions of the adjoint with ratios
in $K$. The least
annihilator is guessed within the bounds of step~4 (at most $32$ unknown
coefficients, entries reconstructed from at most $512$ evaluation points, and
rational coefficients of height about $2^{30}$, lifted from one prime), so a
reduction of high order or degree is missed. All these bounds affect only
completeness: a skipped search never produces a verdict, and a guessed
annihilator is used only after it is proved.

\paragraph{Status.} The module was run with Julia~1.12.7, Nemo~0.56.1 and
AbstractAlgebra~0.50.4. The $53$ tests pass in about five minutes, most of it
for Glaisher's identity and the sum side. The $79$ differences of
Section~\ref{sec:families} take three and a half minutes in all, the $813$ of the
first set of Section~\ref{sec:automata} $17$ minutes, and the $3237$ of the
second set $138$ minutes on one core, $94$ of them for the $29$ that are not
proved. Three runs of \texttt{stress\_quantumwz.jl} with $40$ random cases
each (automata with $2$ to $4$ states) gave no wrong verdict and no exception:
of the $120$ random differences, $113$ were refuted and $7$ proved (equal
marginals), and all $310$ identities were proved, by the zero test or by a
rational certificate without pole, the slowest in $27$ seconds.
\textsc{Oscar} itself could not be installed where this version was prepared.
There, its two-generator PBW algebras were replaced by a small stand-in with
the same interface: normal-ordered terms through \texttt{coefficients} and
\texttt{exponent\_vectors}, the relation $g_2g_1=t\,g_1g_2$, and scalars
accepted only from $K$ or~$\mathbb Z$. An exact reference implementation in
Python, on Flint, was used to cross-check the series, annihilator and channel
stages.

\fussy
\section{The decision question}\label{sec:decision}
Four assertions should be kept apart.
\begin{enumerate}
\item The formal boundary always lies in $\Bhol$ with an effective description
  (Theorem~\ref{thm:holonomic}).
\item Whether the marginal vanishes identically, $C\equiv0$, is decidable
  (Proposition~\ref{prop:zero-exists}). This settles the comparison of
  realizations, but $S(\Delta)=0$ is weaker than $C\equiv0$.
\item Existence of rational certificates, with or without a pole at $T=0$, is
  decidable (Theorem~\ref{thm:decidable}); those of the form $b=P\cdot C$
  correspond to the rational solutions of~\eqref{eq:scalar-cert} for the least
  annihilator of $C$, which the implementation guesses and then proves
  (Proposition~\ref{prop:factor}). Without a pole they are complete
  proofs, but they are not necessary for $S(\Delta)=0$: for identities between
  infinite products every certificate has a pole at $T=0$
  (Proposition~\ref{prop:forced-pole}).
\item Proposed convergence certificates (contraction, channels) are checkable by
  finite exact computation, and enumerating them is a semi-decision
  procedure for $S(\Delta)=0$ \emph{within} those classes. The test of
  Section~\ref{sec:intrinsic} depends only on the sequence $\beta$, and with the
  readings of Section~\ref{sec:readings} it may use the partial sums of any
  reading of each block. It succeeds whenever, for some combination of
  readings, the marginal vanishes or some section of $\beta$ is a sum of
  $q$-hypergeometric terms with ratios in $K$ whose stationary parts cancel in
  pairs through dilation witnesses (within the search bounds); with the
  remainders of Section~\ref{sec:remainders}, it suffices that the partial sums
  of each block are such a sum up to a regular combination of themselves that
  tends to $0$. Completeness beyond this class is not established; the theta
  series of Section~\ref{sec:automata} lies outside it.
\end{enumerate}
A complete, effectively checkable certificate language for convergence would
decide $S(\Delta)=0$: enumerate certificates while computing coefficients of
$S(\Delta)$ until either a certificate verifies or a nonzero coefficient
appears. Equivalently, the problem is decidable iff there is a computable
$N(\Delta)$ such that $t^{N(\Delta)}\mid S(\Delta)$ implies $S(\Delta)=0$.
Neither decidability nor undecidability is established here.

\appendix
\section{Holonomic series do not form a field}\label{app:not-field}
Let $E(Y)=\prod_{j\ge1}(1-t^jY)\in\Bhat\cap K[[Y]]$. Then $E(tY)=E(Y)/(1-tY)$,
$E(0)=1$, and $(1-t^n)[Y^n]E=-t^n[Y^{n-1}]E$, so the coefficients lie in $K$.
Put $g=1+E\in\Bhol\cap\Bhat$; indeed $\bigl((T-t)((1-tY)T-1)\bigr)\cdot g=0$.
Since $g\equiv2\pmod t$, $g^{-1}\in\Bhat$.

\emph{$E$ is transcendental over $K(Y)$.} Otherwise $K(Y,E)$ is a finite
extension of degree $r\ge1$ stable under $\sigma=T$, since
$\sigma(E)=E/(1-tY)$ and $\sigma^{-1}(E)=(1-Y)E$. The norm
$a=N_{K(Y,E)/K(Y)}(E)\in K(Y)^\times$ then satisfies
$a(tY)/a(Y)=(1-tY)^{-r}$. The left side has equal numerator and denominator
degree, the right side does not: a contradiction.

\emph{$g^{-1}$ is not holonomic.} With $P_n=\prod_{j=1}^n(1-t^jY)$ we have
$T^n\cdot g^{-1}=P_n/(P_n+E)$. As rational functions of the transcendental
$E$ over $K(Y)$ these have pairwise distinct simple poles $E=-P_n$, hence are
linearly independent over $K(Y)$. Thus
\[
  g\in\Bhol\cap\Bhat,\qquad g^{-1}\in(\Fhol\cap\Bhat)\setminus\Bhol.
\]

\begin{thebibliography}{10}
\bibitem{GL} S.~Garoufalidis and T.~T.~Q.~L\^e,
\emph{A survey of $q$-holonomic functions}, Enseign. Math. 62 (2016), 501--525.
\href{https://arxiv.org/abs/1601.07487}{arXiv:1601.07487}.
\bibitem{Middeke} J.~Middeke,
\emph{Denominator bounds and polynomial solutions for systems of $q$-recurrences
over $K(t)$ for constant $K$}, ISSAC 2017.
\href{https://arxiv.org/abs/1705.04188}{arXiv:1705.04188}.
\bibitem{DLMF} NIST Digital Library of Mathematical Functions,
\emph{Integer partitions: other restrictions}, \S26.10.
\href{https://dlmf.nist.gov/26.10}{dlmf.nist.gov/26.10}.
\bibitem{Abramov} S.~A.~Abramov, \emph{Rational solutions of linear difference
and $q$-difference equations with polynomial coefficients}, Proc.\ ISSAC'95, ACM,
1995, 285--289.
\bibitem{APP} S.~A.~Abramov, P.~Paule and M.~Petkov\v{s}ek,
\emph{$q$-Hypergeometric solutions of $q$-difference equations},
Discrete Math.\ 180 (1998), 3--22.
\bibitem{PWZ} M.~Petkov\v{s}ek, H.~S.~Wilf and D.~Zeilberger, \emph{A=B},
A K Peters, 1996.
\bibitem{OscarPBW} OSCAR documentation, \emph{Creating PBW-Algebras},
\url{https://docs.oscar-system.org/stable/NoncommutativeAlgebra/PBWAlgebras/creation/}.
\bibitem{OscarSolve} OSCAR documentation, \emph{Solving Linear Systems},
\url{https://docs.oscar-system.org/stable/AbstractAlgebra/how_to_linear_solve/}.
\end{thebibliography}
\end{document}
